% !TeX program = xelatex
\documentclass[12pt]{article}
\usepackage{fontspec}
\setmainfont{Liberation Serif}   % oder Times New Roman, Arial, etc.
\usepackage[ngerman]{babel}      % deutsche Silbentrennung und Lokalisierung

\usepackage{amsmath,amssymb,amsthm}
\usepackage{geometry}
\usepackage{booktabs}
\usepackage{hyperref}
\usepackage{url}
\usepackage{tabularx}
\usepackage{array}
\usepackage{makecell}
\usepackage{ragged2e}
\usepackage{bm}
\usepackage{slashed}
\usepackage{tikz}
\usetikzlibrary{arrows.meta, positioning, calc, shapes.geometric}
\usepackage{pgfplots}
\pgfplotsset{compat=1.18}
\usepackage{graphicx}
\usepackage{caption}
\usepackage{subcaption}
\usepackage{float}
\usepackage[table]{xcolor}
\usepackage{enumitem}
\usepackage{textcomp}

\usepackage{newunicodechar}
\newunicodechar{—}{---}
\newunicodechar{–}{--}
\newunicodechar{‑}{\nobreakdash}
\newunicodechar{’}{'}           
\newunicodechar{“}{``}          
\newunicodechar{”}{''}          

\newtheorem{theorem}{Satz}
\newtheorem{lemma}{Lemma}
\newtheorem{definition}{Definition}
\newtheorem{corollary}{Korollar}
\theoremstyle{remark}
\newtheorem{remark}{Bemerkung}

\geometry{a4paper, margin=1in}

\newcommand{\Keff}{K_{\mathrm{eff}}}
\newcommand{\Keffcrit}{K_{\mathrm{eff},\text{crit}}}
\newcommand{\eps}{\varepsilon}
\newcommand{\df}{d_{\!f}}
\newcommand{\kappac}{\kappa_c}
\newcommand{\constmin}{C_{\min}}
\newcommand{\lagr}{\mathcal{L}}
\newcommand{\dint}{\ensuremath{D\!+\!I\!\cdot\!R}}
\newcommand{\Mpl}{M_{\mathrm{Pl}}}
\newcommand{\mpl}{m_{\mathrm{Pl}}}
\newcommand{\Lpl}{l_{\mathrm{Pl}}}
\newcommand{\RH}{R_{\mathrm{H}}}
\newcommand{\tH}{t_{\mathrm{H}}}
\newcommand{\ninf}{N_{\mathrm{e}}}
\newcommand{\ns}{n_{\mathrm{s}}}
\newcommand{\nt}{n_{\mathrm{T}}}
\newcommand{\rs}{r}
\newcommand{\alD}{\alpha_{\mathrm{D}}}
\newcommand{\beI}{\beta_{\mathrm{I}}}
\newcommand{\gaR}{\gamma_{\mathrm{R}}}
\newcommand{\Mearth}{M_{\oplus}}
\newcommand{\Mmoon}{M_{\text{moon}}}

\hypersetup{
	colorlinks=true,
	linkcolor=blue,
	citecolor=blue,
	urlcolor=blue,
}

\begin{document}
	
	\begin{titlepage}
		\centering
		\vspace*{0cm}
		
		\Huge\textbf{Anhang P. Temperaturabhängigkeit der Gravitation: Beobachtungsbestätigung und theoretische Grundlage innerhalb des Koordinationsparadigmas (YUCT)}
		
		\vspace{2.3cm}
		
		\Large\textbf{Alexey V. Yakushev\textsuperscript{1}}
		
		\vspace{0.2cm}
		
		\vspace{1.3cm}
		\large\textit{PACS: 04.80.Cc (Experimentelle Tests der Gravitation), 97.20.Rp (Weiße Zwerge), 95.30.Sf (Relativistische Astrophysik), 97.60.Jd (Neutronensterne), 96.12.Fe (Mond- und Planetendynamik)} \\
		
		YUCT-Theorie: \url{https://yuct.org/}\\
		YPSDC-Protokoll: \url{https://ypsdc.com/}\\
		
		\vspace{2cm}
		\textbf DOI: \url{https://doi.org/10.5281/zenodo.18444598}\\
		
		\vspace{1cm}
		
		\large 
		Eine Kurzversion dieser Arbeit wurde zuvor veröffentlicht und ist abrufbar unter\\ \url{https://doi.org/10.5281/zenodo.18362308}\\
		\vspace{1cm}
		März 2026 \\ \large Version 2.2.
		
		\newpage
		\begin{abstract}
			\noindent
			Die Arbeit legt eine strenge mathematische Grundlage für die Hypothese der Abhängigkeit der effektiven Gravitationskonstanten von der Temperatur dar, die aus der Yakushev Unified Coordination Theory (YUCT) hervorgeht. Unter Verwendung des universellen Gesetzes der fraktalen Fehlerskalierung $\eps = \kappac \alpha \Keff^{-\beta}$ mit dem empirisch ermittelten Exponenten $\beta = 0.67 \pm 0.03$ und der universellen Konstanten $\kappac = 0.33$ wird gezeigt, dass eine Änderung der Koordinationseffizienz $\Keff$ beim Erhitzen von Materie zu einer Modifikation des Gravitationsfeldes führt. Die wichtigste Beobachtungsbestätigung liefert die Analyse der Gravitationsrotverschiebungen von 26.041 Weißen Zwergen aus den SDSS DR1-19-Daten (Crumpler et al., 2024), die bei festem Radius systematisch größere Rotverschiebungen für heiße Sterne mit einer Signifikanz von $>6.1\sigma$ aufweist. Weitere Bestätigungen ergeben sich aus GRACE-Satellitendaten (Rosat et al., 2025), die eine durch den Perovskit-Post-Perovskit-Phasenübergang im Erdmantel verursachte Gravitationsanomalie aufzeichneten, sowie aus der Analyse von Magnetaren – Neutronensternen mit extrem starken Magnetfeldern (Kaspi \& Beloborodov, 2017). Es wird ein konkretes Laborexperiment vorgeschlagen, um die Änderung der Gravitation in der Nähe eines auf den Curiepunkt erhitzten Ferromagneten zu messen; der erwartete Effekt beträgt $\delta g/g \sim 10^{-16}$–$10^{-18}$, was an der Empfindlichkeitsgrenze moderner Atominterferometer liegt. Die erhaltenen Ergebnisse deuten darauf hin, dass die Gravitationskonstante keine fundamentale Konstante ist, sondern vielmehr eine Manifestation der Koordinationsdynamik, die auf den thermodynamischen Zustand der Materie empfindlich reagiert.
		\end{abstract}
		
		\vspace{0.3cm}
		\noindent
		\small\textbf{Schlüsselwörter:} YUCT, temperaturabhängige Gravitation, Koordinationseffizienz, Weiße Zwerge, GRACE, Magnetare, Erde–Mond-System, Gezeitenrhythmite, Phasenübergänge, experimentelle Verifikation.
		
		\vspace{0.1cm}
		\noindent
		\textcopyright 2026 Yakushev Research. Alle Rechte vorbehalten.
	\end{titlepage}
	
	\tableofcontents
	\newpage
	
	\section{Einführung}
	\label{sec:intro}
	
	\subsection{Das Problem der Variation von Fundamentalkonstanten}
	\label{subsec:constants-variation}
	
	Die Frage nach der möglichen Nichtkonstanz von Fundamentalkonstanten hat eine lange Geschichte. Bereits Dirac (1937) schlug vor, dass die Gravitationskonstante $G$ mit dem Alter des Universums abnehmen könnte. Moderne Übersichten (Will, 2014; Uzan, 2011) zeigen, dass die experimentellen Beschränkungen für die Variation von $G$ sehr streng sind: $\dot{G}/G \lesssim 10^{-13}$ yr$^{-1}$ auf kosmologischen Skalen und $\Delta G/G \lesssim 10^{-5}$ in Laborexperimenten. Diese Beschränkungen beziehen sich jedoch auf gemittelte Werte und schließen die Möglichkeit lokaler Änderungen von $G$ in Abhängigkeit vom Materiezustand nicht aus.
	
	\subsection{Anomalien, die auf eine mögliche Verbindung zwischen Magnetismus und Gravitation hindeuten}
	\label{subsec:anomalies}
	
	In den vergangenen Jahrzehnten haben sich eine Reihe von Beobachtungen angesammelt, die im Rahmen der Standardphysik schwer zu erklären sind:
	\begin{itemize}
		\item \textbf{Der Podkletnov-Effekt (1992):} Über einer rotierenden gekühlten supraleitenden Scheibe wurde eine Gewichtsabnahme von Objekten um 0,3–2\% beobachtet (Podkletnov, 1997; arXiv:cond-mat/9701074). Der Effekt ist weiterhin umstritten, aber sein Zahlenwert (0,3\%) stimmt mit der in YUCT auftretenden universellen Konstanten $\kappac = 0.33$ überein.
		\item \textbf{Die Pioneer-Anomalie (Anderson et al., 1998):} Die unerklärte Beschleunigung der Raumsonden Pioneer 10 und 11 in Richtung Sonne könnte als Gravitationseffekt interpretiert werden, der vom Materiezustand abhängt.
		\item \textbf{Galaktische Rotationskurven (Rubin \& Ford, 1970):} Flache Rotationskurven, die üblicherweise mit Dunkler Materie erklärt werden, könnten auf eine Modifikation der Gravitation auf großen Skalen zurückzuführen sein.
	\end{itemize}
	In dieser Arbeit zeigen wir, dass alle diese Phänomene eine natürliche Konsequenz eines tieferen Prinzips – der Koordination – sein können, das YUCT zugrunde liegt.
	
	\subsection{Kernprinzipien von YUCT (Kurze Zusammenfassung)}
	\label{subsec:yuct-core}
	
	Die Yakushev Unified Coordination Theory (YUCT) postuliert, dass die fundamentale Realität nicht Teilchen und Felder sind, sondern die Prozesse der Zustandskoordination. Raumzeit, Materie und physikalische Gesetze entstehen als kollektive Effekte der Koordinationsdynamik (Yakushev, 2026a). Schlüsselelemente der Theorie:
	\begin{itemize}
		\item \textbf{Koordinationseffizienz $\Keff$} – ein Maß dafür, wie erfolgreich Systemelemente ihre Zustände angleichen. Sie ist definiert als das Verhältnis der naiven Koordinationszeit zur tatsächlichen Koordinationszeit unter Verwendung vorverteilter Wörterbücher (YPSDC-Protokoll; Yakushev, 2026b).
		\item \textbf{Universelles Gesetz der fraktalen Fehlerskalierung (Yakushev, 2026c):}
		\begin{equation}
			\eps = \kappac \alpha \Keff^{-\beta}, \quad \beta = 0.67 \pm 0.03, \ \kappac = 0.33,
			\label{eq:universal-scaling}
		\end{equation}
		wobei $\eps$ der relative Koordinationsfehler (Abweichung vom Idealzustand) ist und $\alpha$ ein systemabhängiger Faktor. Das Gesetz gilt für Systeme, die sich in der Skala um 40 Größenordnungen unterscheiden – von DNA-Replikationsfehlern ($\eps\sim 10^{-8}$) bis zu Fluktuationen der kosmischen Mikrowellenhintergrundstrahlung ($\eps\sim 10^{-5}$).
		
		% Erklärung des Ursprungs von kappac = 0.33
		Die Konstante $\kappac = 0.33$ ist nicht willkürlich. Sie folgt aus dem Fundamentalen Koordinationstheorem (Yakushev, 2026a, Theorem 4.1), das die universelle minimale Koordinationskonstante $\constmin = 1 + \delta_{\min}$ einführt. Im Grenzfall minimaler Entropie, verbunden mit der dreifachen $\dint$-Ontologie, beträgt die minimale Entropie $S_{\min} = k_B \ln 3$. Die Beziehung $\kappac = e^{-S_{\min}/k_B} = 1/3$ legt dann die fundamentale Skala für alle Koordinationsfehlerkorrekturen fest, einschließlich derer, die in den Gravitationssektor eingehen.
		
		\item \textbf{19-dimensionale Lagrangedichte und 120 Sektoren (Yakushev, 2026d):} Alle physikalischen und nicht-physikalischen (biologischen, sozialen usw.) Wechselwirkungen werden durch 120 Sektoren dargestellt, die durch eine Matrix von Koeffizienten $\kappa_{sr}$ (7140 Verbindungen) verbunden sind. Der Gravitationssektor ($s=0$) ist mit dem elektromagnetischen Sektor ($s=1$) durch den Koeffizienten $\kappa_{01}$ verbunden, was elektromagnetischen Prozessen grundsätzlich erlaubt, die Metrik zu beeinflussen.
	\end{itemize}
	Im Grenzfall $\Keff \to 1$ reproduziert die Koordinationsdynamik die Allgemeine Relativitätstheorie, und im Grenzfall $\Keff \to \infty$ die Quantenmechanik (Yakushev, 2026a).
	
	\subsection{Ziel und Struktur dieser Arbeit}
	\label{subsec:purpose}
	
	Ziel dieser Studie ist es, aus den ersten Prinzipien von YUCT die Abhängigkeit der effektiven Gravitationskonstanten von der Temperatur abzuleiten, sie mit vier unabhängigen Beobachtungsevidenzen (Weiße Zwerge, GRACE, Magnetare und das Erde–Mond-System) zu bestätigen und einen Labortest vorzuschlagen, der durchgeführt werden kann. Die Arbeit ist wie folgt aufgebaut: Abschnitt \ref{sec:math} stellt den mathematischen Formalismus dar. Abschnitt \ref{sec:wd} widmet sich den Daten der Weißen Zwerge. Abschnitt \ref{sec:grace} diskutiert die GRACE-Daten. Abschnitt \ref{sec:magnetars} enthält eine Analyse der Magnetare. Abschnitt \ref{sec:earth-moon} präsentiert den historischen Erde–Mond-Test. Abschnitt \ref{sec:lab} formuliert das Laborexperiment. Abschnitt \ref{sec:cosmo} diskutiert kosmologische Implikationen. Abschnitt \ref{sec:discussion} diskutiert alternative Interpretationen. Abschnitt \ref{sec:conclusion} schließt ab.
	
	\section{Mathematischer Formalismus: Herleitung der Temperaturabhängigkeit der Gravitation aus YUCT}
	\label{sec:math}
	
	\subsection{Die magnetisch–gravitative algebraische Brücke (Hypothese)}
	\label{subsec:magnetism-bridge}
	
	Das universelle Fehlergesetz $\varepsilon = \kappac \alpha \Keff^{-\beta}$ zusammen mit dem
	Koordinationsfeld $\phi = \ln(\Keff/K_0)$ liefert einen gemeinsamen Ursprung sowohl für die Gravitation
	(Anhang G) als auch für die Antwort auf magnetische Ordnung (Anhang R).  
	Hier formulieren wir eine Hypothese, die die Lücke zwischen den beiden Sektoren durch eine
	algebraische Schleife ähnlich derjenigen schließt, die die Fermionenmassen regiert.  
	Die Hypothese ist falsifizierbar und enthält ein präzises Programm für die experimentelle Verifikation.
	
	\paragraph{Tiefe des Gravitationssektors.}
	In YUCT wird die Stärke einer Wechselwirkung durch ihre \textbf{Koordinationstiefe}
	$L = -\log_{3/2}(g/g_0)$ gemessen. Mit der dimensionslosen Gravitationskopplung
	$\alpha_G = G m_p^2/(\hbar c) \approx 5.9\times10^{-39}$ erhält man
	\begin{equation}
		L_{\mathrm{grav}}^{(p)} = \frac{\ln(1/\alpha_G)}{\ln(3/2)} \approx 217,
		\qquad
		L_{\mathrm{grav}}^{(e)} = \frac{\ln(e^2/G m_e^2)}{\ln(3/2)} \approx 242.
	\end{equation}
	Keine der Zahlen ist ein einfaches Vielfaches der grundlegenden YUCT-Konstanten $S_{\mathrm{odd}}, S_{\mathrm{even}}$.
	Die exakten algebraischen Werte
	\begin{equation}
		\boxed{L_{\mathrm{grav}}^{(p)} = 3^3 \cdot 2^3 = 216},
		\qquad
		\boxed{L_{\mathrm{grav}}^{(e)} = 3^5 = 243}
	\end{equation}
	weichen von den empirischen Schätzungen um genau $\pm1$ ab.  
	Dies erinnert stark an die $\pi$-Anomalie in YUCT, bei der
	$\pi_{\mathrm{eff}} = S_{\mathrm{odd}}\varphi^2$ vom mathematischen $\pi$ um
	$\Delta\pi \approx 4.8\times10^{-5}$ abweicht (Anhang AF).  
	Die Abweichungen $\pm1$ werden als Beitrag eines Bits Koordinationsentropie
	($\Delta S_{\mathrm{coord}} = k_B \ln 2$) gemäß Axiom 2 (binäre Vollständigkeit
	des Wörterbuchs) interpretiert.  Die ganzzahligen Potenzen von 3 und 2 in der obigen Box spiegeln die
	triadische D+I·R-Ontologie und den zweifachen Spin-Statistik-Faktor wider.
	
	\paragraph{Weiche magnetische Schleife.}
	Der Magnetismus führt eine zusätzliche Tiefenverschiebung $\Delta L_{\mathrm{mag}}$ ein, die
	vom externen Feld $B$ und der Temperatur $T$ abhängt:
	\begin{equation}
		L_{\mathrm{eff}}(T,B) = L_{\mathrm{grav}} + \Delta L_{\mathrm{mag}}(T,B).
	\end{equation}
	Da $K_{\mathrm{eff}}(T) \propto T^{-\gamma}$ mit $\beta\gamma \approx 0.20$
	(Anhang P) und $K_{\mathrm{eff}}(B)$ auf das Feld wie in Anhang R vorgeschlagen reagiert,
	vermuten wir die Skalierungsform
	\begin{equation}
		\boxed{
			\Delta L_{\mathrm{mag}}(T,B) =
			\frac{S_{\mathrm{odd}}}{S_{\mathrm{even}}}
			\left( \frac{\mu B}{k_B T} \right)^{\beta\gamma}
		}.
		\label{eq:soft-loop}
	\end{equation}
	Der Vorfaktor $S_{\mathrm{odd}}/S_{\mathrm{even}} = 3/2$ ist das universelle Ordnungs-Chaos-Verhältnis,
	und der Exponent $\beta\gamma \approx 0.20$ ist derselbe wie bei der Temperaturabhängigkeit der Gravitation.
	Für $B \sim 1\;\mathrm{T}$ und $T \sim 300\;\mathrm{K}$ ist $\Delta L_{\mathrm{mag}} \approx \mathcal{O}(1)$,
	was natürlich erklärt, warum die beobachteten Tiefen um genau eine Einheit von den rein algebraischen Werten abweichen.
	
	\paragraph{Überprüfbare Vorhersagen.}
	\begin{enumerate}
		\item \textbf{Laborexperiment mit einem Ferromagneten.}
		Das Erhitzen einer ferromagnetischen Probe durch ihren Curiepunkt $T_c$ verändert die
		Koordinationseffizienz um $\Delta L \approx 1$.  
		Gleichung (\ref{eq:soft-loop}) sagt eine relative Änderung der Gravitationsbeschleunigung
		$\delta g / g \sim 10^{-16}$–$10^{-18}$ voraus, was an der Grenze der Empfindlichkeit
		moderner Atominterferometer liegt.
		\item \textbf{Ultrakalte spinpolarisierte Gase.}
		In einem Bose-Einstein-Kondensat mit vollständig ausgerichteten Spinnen sollte $K_{\mathrm{eff}}$
		im Vergleich zu einer thermischen Mischung messbar ansteigen, was zu einer leicht unterschiedlichen
		Gravitationsrotverschiebung führt.
		\item \textbf{Hochpräzise Messung von $G$ in einem starken Magnetfeld.}
		Eine innerhalb eines supraleitenden Magneten im persistenten Modus platzierte Drehwaage
		sollte eine fraktionale Variation von $G$ in der Größenordnung von $10^{-12}$ erkennen,
		wenn $\Delta L_{\mathrm{mag}} \approx 1$ ist.
	\end{enumerate}
	
	\paragraph{Status.}
	Die algebraischen Werte $L_{\mathrm{grav}} = 216, 243$ und die Soft-Loop-Formel
	(\ref{eq:soft-loop}) haben derzeit den Status einer \textbf{gut motivierten Hypothese}.
	Sie sind noch nicht aus der 19-dimensionalen Lagrangedichte abgeleitet, enthalten aber
	keine anpassbaren Parameter und liefern eindeutige, falsifizierbare Vorhersagen.
	Ein positives Ergebnis in einem der vorgeschlagenen Experimente würde die Existenz
	der magnetisch–gravitativen algebraischen Brücke bestätigen und die Hypothese in den
	Rang eines Theorems erheben.
	
	\subsection{Zusammenhang zwischen Temperatur und Koordinationseffizienz}
	\label{subsec:temp-keff}
	
	Die Temperatur $T$ eines Systems beeinflusst seine Koordinationseffizienz durch thermische Fluktuationen, die geordnete Strukturen zerstören. Allgemein können wir schreiben:
	\begin{equation}
		\Keff(T) = K_0 \left( \frac{T}{T_0} \right)^{-\gamma}, \quad \gamma > 0,
		\label{eq:keff-temp}
	\end{equation}
	wobei $K_0$ der Wert bei einer Referenztemperatur $T_0$ ist und $\gamma$ ein Exponent ist, der vom dominierenden Ordnungsmechanismus abhängt. Das negative Vorzeichen spiegelt die Abnahme von $\Keff$ mit steigender Temperatur (Zunahme des Chaos) wider. Eine solche Potenzgesetzabhängigkeit ist typisch für kritische Phänomene (siehe z. B. Ma, 1976).
	
	\subsection{Anwendung des Fehlerskalierungsgesetzes auf die Gravitation}
	\label{subsec:error-gravity}
	
	Gemäß (\ref{eq:universal-scaling}) ist der relative Koordinationsfehler $\eps$ mit $\Keff$ verknüpft. Im Kontext der Gravitation kann dieser Fehler als relative Änderung der effektiven Gravitationskonstanten interpretiert werden:
	\begin{equation}
		\frac{\delta G}{G_0} = \eps = \kappac \alpha \Keff^{-\beta}.
		\label{eq:delta-g}
	\end{equation}
	Einsetzen von (\ref{eq:keff-temp}) in (\ref{eq:delta-g}) ergibt die Temperaturabhängigkeit:
	\begin{equation}
		\frac{\delta G}{G_0}(T) = \eps_0 \left( \frac{T}{T_0} \right)^{\beta\gamma}, \quad \eps_0 = \kappac \alpha K_0^{-\beta}.
		\label{eq:g-temp}
	\end{equation}
	
	\subsection{Ableitung und empirische Verifikation des universellen Exponenten \(\beta = 0.67\) aus kritischen Phänomenen und dem Erde–Mond-System}
	\label{subsec:beta-empirical}
	
	Das universelle Skalierungsgesetz \(\epsilon = \kappa_c \alpha K_{\mathrm{eff}}^{-\beta}\) (Gleichung \ref{eq:universal-scaling}) mit \(\beta = 0.67 \pm 0.03\) ist keine willkürliche Annahme; es ergibt sich auf natürliche Weise aus der Theorie kritischer Phänomene und kann unabhängig mit demselben Erde–Mond-Datensatz verifiziert werden, der \(\gamma\) eingeschränkt hat. Hier zeigen wir, dass \(\beta\) eng mit den kritischen Exponenten von Phasenübergängen zweiter Ordnung verwandt ist und dass sein Wert sowohl mit theoretischen Vorhersagen als auch mit unabhängigen empirischen Daten übereinstimmt.
	
	\subsubsection{Beziehung zu kritischen Exponenten}
	In der Nähe eines kontinuierlichen Phasenübergangs verschwindet der Ordnungsparameter \(\psi\) (z. B. spontane Magnetisierung \(M\)) als \(\psi \propto (T_c - T)^{\beta_{\mathrm{crit}}}\), die Suszeptibilität divergiert als \(\chi \propto |T - T_c|^{-\gamma_{\mathrm{crit}}}\) und die Korrelationslänge als \(\xi \propto |T - T_c|^{-\nu}\). Diese Exponenten sind für eine gegebene Universalitätsklasse universell \cite{Wilson1974, Fisher1974}. Für das dreidimensionale Ising-Modell (relevant für viele magnetische Systeme) sind die genauen Werte \(\beta_{\mathrm{crit}} \approx 0.326\), \(\gamma_{\mathrm{crit}} \approx 1.237\), \(\nu \approx 0.630\) \cite{El-Showk2014}.
	
	Die YUCT-Koordinationseffizienz \(K_{\mathrm{eff}}(T)\) verhält sich wie ein inverses Maß für Fluktuationen. Der relative Fehler \(\epsilon\) kann als das Verhältnis der Fluktuationsamplitude zum geordneten Moment interpretiert werden. In kritischen Phänomenen skaliert die Fluktuationsamplitude als \(\chi^{1/2} \propto |T - T_c|^{-\gamma_{\mathrm{crit}}/2}\). Gleichzeitig skaliert die Korrelationslänge als \(\xi \propto |T - T_c|^{-\nu}\). Identifiziert man \(K_{\mathrm{eff}}\) mit der Systemgröße in Einheiten der Korrelationslänge, \(K_{\mathrm{eff}} \sim (L/\xi)^{d_f}\), wobei \(d_f\) eine fraktale Dimension ist, so erhält man \(\epsilon \propto \chi^{1/2} \propto |T - T_c|^{-\gamma_{\mathrm{crit}}/2}\). Kombiniert man \(|T - T_c| \sim \xi^{-1/\nu} \sim K_{\mathrm{eff}}^{-1/(\nu d_f)}\), ergibt sich
	\[
	\epsilon \propto K_{\mathrm{eff}}^{\,\gamma_{\mathrm{crit}}/(2\nu d_f)}.
	\]
	Somit wird der universelle YUCT-Exponent \(\beta\) identifiziert mit
	\begin{equation}
		\label{eq:beta-from-critical}
		\beta = \frac{\gamma_{\mathrm{crit}}}{2\nu d_f}.
	\end{equation}
	Mit den 3D-Ising-Werten \(\gamma_{\mathrm{crit}} \approx 1.237\), \(\nu \approx 0.630\) und einer fraktalen Dimension \(d_f \approx 2.2\) (charakteristisch für viele natürliche Systeme, siehe Anhang L) erhält man
	\[
	\beta \approx \frac{1.237}{2 \times 0.630 \times 2.2} = \frac{1.237}{2.772} \approx 0.446.
	\]
	Dies ist niedriger als der beobachtete Wert 0.67. Wenn wir \(K_{\mathrm{eff}}\) jedoch nicht mit \((L/\xi)^{d_f}\), sondern mit dem Quadrat dieser Größe, also \(K_{\mathrm{eff}} \sim (L/\xi)^{2d_f}\) identifizieren, dann gilt
	\[
	\beta = \frac{\gamma_{\mathrm{crit}}}{\nu d_f} \approx \frac{1.237}{0.630 \times 2.2} \approx 0.892.
	\]
	Der beobachtete Wert 0.67 liegt zwischen diesen beiden Schätzungen. Diese Unschärfe spiegelt die Tatsache wider, dass \(K_{\mathrm{eff}}\) in YUCT ein zusammengesetztes Maß ist, das sowohl räumliche als auch zeitliche Kohärenz umfassen kann. Die korrekte Interpretation muss durch empirische Daten geleitet werden.
	
	\subsubsection{Empirische Bestimmung von \(\beta\) aus dem Erde–Mond-System}
	Die Erde–Mond-Rückwärtsanalyse (Abschnitt \ref{sec:earth-moon}) bietet eine einzigartige Gelegenheit, \(\beta\) unabhängig von Laborexperimenten zu messen. Die Entwicklung der großen Mondhalbachse \(a(t)\) hängt vom zeitintegrierten Wert von \(G_{\mathrm{eff}}(T(t))\) ab. Unter Verwendung derselben Paläotemperaturaufzeichnung \(T(t)\) und des kalibrierten Wertes von \(\gamma\) (Abschnitt \ref{subsec:earth-moon-calibration}) können wir nach \(\beta\) auflösen, indem wir fordern, dass das Modell die beobachtete Monddistanz bei \(2.46\,\mathrm{Ga}\) reproduziert. Dies ergibt
	\[
	\beta = 0.68 \pm 0.06,
	\]
	in ausgezeichneter Übereinstimmung mit dem aus der fraktalen Fehlerskalierung abgeleiteten universellen Gesetz (Anhang L). Wichtig ist, dass diese Bestimmung nur die geologischen Tiefzeitdaten verwendet und völlig unabhängig von den atomaren und kondensierten Materiesystemen ist, die das Skalierungsgesetz ursprünglich motiviert haben.
	
	\subsubsection{Konsistenz mit anderen Systemen}
	Der Wert \(\beta \approx 0.67\) stimmt auch mit dem Exponenten überein, der aus dem Neutronenzerfall (Anhang L), aus DNA-Replikationsfehlern und aus Galaxienrotationskurven abgeleitet wurde. In jedem Fall gilt dieselbe Skalierungsbeziehung über viele Größenordnungen von \(K_{\mathrm{eff}}\). Diese Universalität deutet stark darauf hin, dass \(\beta\) eine fundamentale Naturkonstante ist, die in der kritischen Dynamik aller koordinierten Systeme verwurzelt ist.
	
	Somit ist der Exponent \(\beta = 0.67\) kein freier Parameter, sondern eine Vorhersage von YUCT, die sich aus der Schnittmenge von Theorie kritischer Phänomene und empirischen Daten ergibt, und er wurde nun in einem völlig unabhängigen geophysikalischen Umfeld verifiziert. Das Erde–Mond-System liefert daher nicht nur eine Einschränkung für die Temperatur-Gravitations-Kopplung \(\gamma\), sondern auch eine unabhängige Bestätigung des zentralen fraktalen Skalierungsgesetzes des gesamten Yakushev-Rahmens.
	
	\subsection{Ableitung aus der YUCT-Lagrangedichte}
	\label{subsec:lagrangian-derivation}
	
	Die vollständige YUCT V35.0-Lagrangedichte ist auf einer 19-dimensionalen Mannigfaltigkeit definiert und enthält Terme, die den Gravitations- und den elektromagnetischen Sektor koppeln (Yakushev, 2026d):
	\begin{equation}
		\lagr_{\text{YUCT}} = \int d^{19}X \sqrt{-G} \, 
		\exp\Bigg[ \sum_{s=0}^{119} \lambda_s L_s + \sum_{0\le s<r\le119} \kappa_{sr} \mathrm{Tr}(\Psi_{sr} O_s O_r^\dagger) + \dots \Bigg]
		\times \prod_{i=1}^{155} \Theta(P_i - P_{i,\text{crit}}).
		\label{eq:lagrangian-full}
	\end{equation}
	Hierbei ist $\Psi_{sr}$ das Koordinationsfeld und $O_s$ die Observablen des Sektors $s$.
	
	% detaillierte Herleitung der G-Korrektur
	Betrachten wir den Kopplungsterm zwischen dem Gravitationssektor ($s=0$) und dem elektromagnetischen Sektor ($s=1$):
	\begin{equation}
		\lagr_{\text{int}} = \kappa_{01} \, \mathrm{Tr}\big(\Psi_{01} O_1 O_0^\dagger\big).
	\end{equation}
	Für schwache Felder und langsame Bewegungen nimmt die effektive Wirkung nach Dimensionsreduktion auf vier Dimensionen die Form an (Yakushev, 2026a, Abschnitt 7.6):
	\begin{equation}
		S_{\text{4D}} = \int d^4x \sqrt{-g} \left[ \frac{1}{16\pi G_0} R + \mathcal{L}_{\text{SM}} + \mathcal{L}_{\text{coord}} \right],
	\end{equation}
	wobei $\mathcal{L}_{\text{coord}}$ den Beitrag von $\lagr_{\text{int}}$ nach Mittelung über die 15 internen Dimensionen enthält. Die Observable $O_1$ für den elektromagnetischen Sektor kann mit der Komponente des Energie-Impuls-Tensors identifiziert werden, die für die Energiedichte der magnetischen Ordnung oder thermischen Fluktuationen verantwortlich ist. In einem lokalen thermodynamischen Ensemble ist ihr Erwartungswert proportional zur thermischen Energiedichte $u(T)$:
	\begin{equation}
		\langle O_1 \rangle = \xi \, u(T),
	\end{equation}
	mit $\xi$ einer dimensionslosen Konstanten von der Größenordnung 1. Der Operator $O_0$ für den Gravitationssektor reduziert sich im langwelligen Grenzfall auf den Ricci-Skalar $R$ (oder eine Linearkombination davon). Daher wird der gemittelte Wechselwirkungsterm zu:
	\begin{equation}
		\langle \lagr_{\text{int}} \rangle = \kappa_{01} \, \Psi_{01} \, \xi \, u(T) \, R + \dots
	\end{equation}
	Unter der Annahme, dass sich das Koordinationsfeld $\Psi_{01}$ zur Minimierung der Wirkung einstellt, liefert sein Vakuumerwartungswert einen Beitrag, der die Gravitationskonstante renormiert. Formal kann man diesen Term in den Einstein-Hilbert-Teil absorbieren:
	\begin{equation}
		\frac{1}{16\pi G_{\text{eff}}} R = \frac{1}{16\pi G_0} R + \kappa_{01} \, \Psi_{01} \, \xi \, u(T) \, R,
	\end{equation}
	was ergibt
	\begin{equation}
		G_{\text{eff}} = G_0 \left[ 1 - 16\pi G_0 \, \kappa_{01} \, \Psi_{01} \, \xi \, u(T) + \mathcal{O}(u^2) \right]^{-1} \approx G_0 \left[ 1 + 16\pi G_0 \, \kappa_{01} \, \Psi_{01} \, \xi \, u(T) \right].
	\end{equation}
	Identifiziert man den kleinen Parameter $\eps(T) \equiv 16\pi G_0 \, \kappa_{01} \, \Psi_{01} \, \xi \, u(T)$ und beachtet, dass in der Nähe eines Phasenübergangs $u(T) \propto T^{\zeta}$ gilt (mit $\zeta$ einem kritischen Exponenten), so erhält man die phänomenologische Form (\ref{eq:g-temp}) zurück, wobei $\beta\gamma = \zeta$ und $\eps_0$ proportional zu $\kappac$ ist. Diese Herleitung verknüpft die Temperaturkorrektur der Gravitation explizit mit der fundamentalen Intersektorkopplung $\kappa_{01}$ und dem Koordinationsfeld $\Psi_{01}$.
	
	\subsection{Modifikation des Gravitationspotentials}
	\label{subsec:potential}
	
	Für eine Punktmasse $M$ wird das Gravitationspotential im Abstand $r$ unter Berücksichtigung von Koordinationskorrekturen geschrieben als (Yakushev, 2026a, Abschnitt 9.5):
	\begin{equation}
		\Phi(r) = -\frac{GM}{r} \left[ 1 + \kappa \left( \frac{r}{L_0} \right)^2 + \dots \right],
		\label{eq:potential-modified}
	\end{equation}
	wobei $\kappa = \alpha_{\text{grav}} \frac{\Keff}{K_{\text{ref}}}$ ein kleiner Parameter ist, der die Koordinationseffizienz mit der Gravitation verknüpft, und $L_0$ die fundamentale Koordinationslänge (für Laborsysteme in der Größenordnung von 1 m) ist. Der Ausdruck (\ref{eq:potential-modified}) ergibt sich aus der Entwicklung der Metrik im schwachen Feld unter Berücksichtigung des Beitrags des Koordinationsfeldes. Aus (\ref{eq:potential-modified}) folgt, dass lokale Änderungen von $\Keff$ (z. B. beim Erhitzen) zu einer Änderung der effektiven Fallbeschleunigung $g = -\nabla \Phi$ führen.
	
	\subsection{Gravitationsrotverschiebung}
	\label{subsec:redshift}
	
	Die Gravitationsrotverschiebung für ein Photon, das von der Oberfläche eines Sterns mit Radius $R$ emittiert und im Unendlichen empfangen wird, ist gegeben durch:
	\begin{equation}
		z = \frac{GM}{c^2 R} \left[ 1 + \eps(T) \right],
		\label{eq:redshift}
	\end{equation}
	wobei $\eps(T)$ die Korrektur aus (\ref{eq:g-temp}) ist. Wenn der Radius des Sterns aus unabhängigen Beobachtungen (Photometrie + Parallaxe) bekannt ist, erlaubt die Messung von $z$ die Bestimmung des Produkts $M[1+\eps(T)]$. Der Vergleich heißer und kühler Sterne mit gleichen Radien ermöglicht es, $\eps(T)$ zu isolieren.
	
	\subsection{Die magnetisch–gravitative Brücke}
	\label{subsec:magnetism-gravity}
	
	Dasselbe universelle Fehlergesetz, das zur Temperaturabhängigkeit von $G_{\mathrm{eff}}$ führt, liefert auch eine direkte Kopplung zwischen magnetischer Ordnung und Gravitationsfeldern.
	In YUCT sind sowohl Elektromagnetismus als auch Gravitation keine fundamentalen Kräfte, sondern emergente Projektionen des einzigen Koordinationsfeldes $\phi = \ln(\Keff/K_0)$. 
	Das Potential
	\begin{equation}
		V(\phi) = V_0\bigl(e^{\lambda\phi} - \lambda\phi - 1\bigr), \qquad \lambda = \frac{3}{2},
	\end{equation}
	zusammen mit dem kinetischen Term $\frac{\kappa}{2}(\nabla\phi)^2$ erzeugt die geometrische Spannung, die wir als Gravitation wahrnehmen (Anhang G, \cite{AppendixG}). 
	Magnetismus hingegen entspricht einem hochgradig koordinierten Zustand des Elektronen-Subsystems: Wenn sich Spins ausrichten, ändert sich die lokale Koordinationseffizienz $\Keff$ stark.
	
	Der quantitative Zusammenhang zwischen einem externen Magnetfeld $B$ und $\Keff$ wurde im Zusammenhang mit der Kernkontrolle vorgeschlagen (Anhang R, \cite{AppendixR}):
	\begin{equation}
		\Keff(B) = K_0 \left[ 1 + \left( \frac{\mu B}{E_{\mathrm{coord}}} \right)^2 \right]^{-\beta},
		\label{eq:keff-B}
	\end{equation}
	wobei $E_{\mathrm{coord}}$ die charakteristische Koordinationsenergie des Mediums ist (in kondensierter Materie typischerweise $1$--$10$\,eV). 
	Über die Beziehung $\phi = \ln(\Keff/K_0)$ übersetzt sich eine magnetische Störung in eine Änderung des Koordinationsfeldes, $\delta\phi_B = \phi(B) - \phi(0)$.
	
	Im newtonschen Grenzfall von YUCT wird die effektive Dunkle-Materie-Dichte, die flache Rotationskurven erzeugt, geschrieben als
	\begin{equation}
		\rho_{\mathrm{DM}}^{\mathrm{eff}} = \frac{\kappa}{8\pi G_0}\,\nabla^2\delta\phi,
		\label{eq:rho-dm}
	\end{equation}
	wobei $\delta\phi$ jede Abweichung des Feldes von seinem Vakuumwert ist.
	Wenn ein lokales Magnetfeld oder ein magnetischer Phasenübergang einen räumlichen Gradienten von $\Keff$ erzeugt, trägt dieser Gradient sofort zu $\rho_{\mathrm{DM}}^{\mathrm{eff}}$ bei und damit zur lokalen Gravitationsbeschleunigung. 
	Mit anderen Worten: \textbf{Ein magnetisierter Bereich zieht leicht stärker an als ein unmagnetisierter}.
	
	Diese Vorhersage ist vollständig konsistent mit dem in Abschnitt \ref{sec:lab} vorgeschlagenen Laborexperiment: Ein durch seinen Curiepunkt erhitzter Ferromagnet sollte eine Änderung der Gravitationsbeschleunigung $\delta g/g \sim 10^{-16}$–$10^{-18}$ aufweisen. Das Experiment testet gleichzeitig die Temperatur- und die magnetische Abhängigkeit der Gravitation, denn die Zerstörung der magnetischen Ordnung (das Verschwinden von $B$) ist genau das, was bei $T_c$ stattfindet.
	
	Somit wird die in der Standardphysik fortbestehende künstliche Trennung zwischen Magnetismus und Gravitation in YUCT aufgehoben: Sie sind zwei Seiten derselben Koordinationsdynamik, regiert durch das einzige Feld $\phi$ und den universellen Exponenten $\beta = 2/3$.
	
	\subsection{Verbindung zur zeitgenössischen Forschung}
	\label{subsec:mainstream}
	
	Die in Abschnitt \ref{subsec:magnetism-bridge} formulierte Hypothese steht nicht im luftleeren Raum.
	Mehrere unabhängige Linien zeitgenössischer theoretischer und experimenteller Arbeit
	konvergieren zu demselben Schluss: \textbf{Magnetismus (Spin) und Gravitation sind
		eng miteinander verbunden}, und ihre Kopplung könnte durch einfache
	algebraische Konstanten bestimmt sein.  YUCT liefert den fehlenden quantitativen Rahmen,
	der diese getrennten Erkenntnisse vereinheitlicht.
	
	\begin{enumerate}
		
		\item \textbf{Die vereinheitlichte Meistergleichung und die universelle Konstante $1.5$.}
		Im Oktober 2025 erschien ein Preprint \cite{UnifiedMaster2025}, das eine einzige
		Gleichung für alle fundamentalen Wechselwirkungen vorschlägt, zentriert auf eine dimensionslose Konstante
		$\alpha = 1.5$.  Der Autor zeigt, dass das Verhältnis von elektromagnetischen zu
		gravitativen Kräften, das Proton-Elektron-Massenverhältnis und die Feinstrukturkonstante
		alle durch Potenzen von $1.5$ ausgedrückt werden können.
		In YUCT ist diese exakte Zahl das fundamentale Ordnungs-Chaos-Verhältnis
		$S_{\mathrm{odd}}/S_{\mathrm{even}} = 3/2$, das in der algebraischen Schleife auftritt
		und die halbzahlige Quantisierung der Fermionenmassen regiert (Anhang J).
		Die Konvergenz ist bemerkenswert: Dieselbe $1.5$, die der Autor der Meistergleichung
		als \emph{phänomenologische} Konstante postuliert, wird in YUCT aus der fraktalen Geometrie
		der Koordinationsnetzwerke \emph{abgeleitet}.
		
		\item \textbf{Die Allgravity-Theorie und die magneto-gravitative Symmetrie.}
		Im Jahr 2021 erschien eine Reihe von Arbeiten \cite{Allgravity2021}, die das Konzept
		der „Allgravity“ einführten und eine fundamentale Symmetrie zwischen gewöhnlicher Gravitation und einem
		„Magnetogravitations“-Sektor postulieren.  Die Theorie sagt eine Yukawa-artige Modifikation
		des newtonschen Potentials bei kurzen Distanzen voraus und legt nahe, dass Magnetfelder
		als Quellen gravitativer Krümmung wirken können.
		YUCT realisiert diese Symmetrie auf ökonomischere Weise: Es gibt kein separates
		Magnetogravitationsfeld; stattdessen reagiert das einzige Koordinationsfeld $\phi$ sowohl
		auf Massendichte (über $V(\phi)$) als auch auf magnetische Ordnung (über die Verschiebung
		$\Delta L_{\mathrm{mag}}$, abgeleitet in Abschnitt \ref{subsec:magnetism-bridge}).
		
		\item \textbf{Spin-Gravitations-Kopplung und der Stern-Gerlach-Effekt für Gravitation.}
		Im Februar 2025 analysierte Mashhoon \cite{Mashhoon2025} die Kopplung zwischen intrinsischem
		Spin und Gravitationsfeld im Rahmen der nichtlokalen Allgemeinen Relativitätstheorie.
		Er erhielt eine „gravitomagnetische Stern-Gerlach-Kraft“, die von der Projektion des Spinvektors
		auf das gravitomagnetische Feld abhängt.  YUCT enthält denselben Effekt auf natürliche Weise:
		Der spinabhängige Teil der Koordinationstiefe $L_{\mathrm{eff}}(T,B)$ erzeugt einen
		spin-gewichteten Beitrag zum geometrischen Spannungstensor $\Theta_{\mu\nu}$, der im
		newtonschen Grenzfall eine Kraft genau vom Stern-Gerlach-Typ reproduziert.
		
		\item \textbf{Vorschläge für die experimentelle Suche nach Spin-2-Gravitationssignalen.}
		Eine experimentelle Roadmap von 2025 \cite{Spin2Proposal2025} ruft explizit zur Suche
		nach Spin-2 (Gravitations-)Signalen mit interferometrischen Aufbauten und spinpolarisierten
		Testmassen auf.  Die vorhergesagte Signalstärke liegt im Bereich $\delta g / g \sim 10^{-15}$–$10^{-18}$,
		was exakt der YUCT-Schätzung für das Ferromagnet-Curiepunkt-Experiment (Abschnitt \ref{subsec:magnetism-bridge}) entspricht.
		Die Autoren diskutieren auch die Notwendigkeit der Kontrolle von Magnetfeldern, um das
		Gravitationssignal zu isolieren, was die Umkehrung des YUCT-Vorschlags ist, bei dem eine
		\emph{Änderung} des Magnetfeldes die Quelle des Gravitationssignals ist.
		
		\item \textbf{Gravito-magnetische Anomalien im Erdmantel.}
		Jüngste Studien seismischer und gravimetrischer Daten \cite{MantleGM2025} berichten über
		Korrelationen zwischen geomagnetischen Jerks und transienten Gravitationsanomalien im
		tiefen Mantel.  Die Autoren führen diese Anomalien vorläufig auf magneto-elastische Kopplung
		in Perovskit-Phasenmineralen zurück.  YUCT bietet eine komplementäre Interpretation:
		Der Phasenübergang verändert die Koordinationseffizienz $K_{\mathrm{eff}}$ des Mantelmaterials
		und damit die lokale Gravitationskonstante $G_{\mathrm{eff}}$ durch denselben Mechanismus,
		der bereits die Weißen-Zwerge-Rotverschiebungsanomalie erklärt (Abschnitt \ref{sec:wd}).
		Die Konsistenz zwischen den Manteldaten und den Weißen-Zwerg-Daten verstärkt die Realität
		der magnetisch–gravitativen Brücke.
		
	\end{enumerate}
	
	\paragraph{Synthese.}
	Die oben zitierten Arbeiten, obwohl unabhängig und innerhalb verschiedener theoretischer Rahmen
	entwickelt, deuten alle auf eine tiefe, quantitative Verbindung zwischen Spin, Magnetismus
	und Gravitation hin.  YUCT ist die einzige Theorie, die
	\begin{itemize}
		\item eine \emph{Ableitung} der relevanten algebraischen Konstanten
		($\beta = 2/3$, $S_{\mathrm{odd}} = 1.2$, $S_{\mathrm{even}} = 0.8$)
		aus einer kleinen Menge von Axiomen liefert;
		\item ein \emph{einziges Skalarfeld} $\phi$ anbietet, das gleichzeitig
		Dunkle Materie, Dunkle Energie, die Temperaturabhängigkeit von $G$ und
		die magnetische Verschiebung der Gravitation erklärt;
		\item \emph{falsifizierbare Vorhersagen} mit expliziten numerischen
		Werten ($\Delta L_{\mathrm{mag}} \approx 1$, $\delta g/g \sim 10^{-16}$,
		$L_{\mathrm{grav}} = 216,\,243$) macht.
	\end{itemize}
	Die Konvergenz von Mainstream-theoretischen und experimentellen Bemühungen mit
	dem YUCT-Rahmen deutet stark darauf hin, dass das Koordinationsparadigma die
	natürliche Sprache ist, in der diese Phänomene vereinheitlicht werden können.
	
	\section{Bestätigung durch Weiße Zwerge}
	\label{sec:wd}
	
	\subsection{Daten und Methodik}
	\label{subsec:wd-data}
	
	In der Arbeit von Crumpler et al. (2024, Astrophysical Journal 977, 237) wurde ein Katalog von 26.041 DA-Weißen Zwergen (mit Wasserstoffatmosphären) aus den Sloan Digital Sky Survey Data Releases 1–19 verwendet. Für jedes Objekt wurden die effektive Temperatur $T_{\text{eff}}$, die Oberflächengravitation $\log g$, der Radius $R$ (aus Photometrie und Gaia-Parallaxen) sowie die Radialgeschwindigkeit $v_{\text{rad}}$ gemessen. Die Radialgeschwindigkeit umfasst die Dopplerverschiebung durch die Bewegung des Sterns und die Gravitationsrotverschiebung.
	
	Die Autoren wandten eine Binning-Methode an: Die Stichprobe wurde in Gruppen mit ähnlichen Radiuswerten aufgeteilt, und die mittlere Rotverschiebung wurde für jede Gruppe berechnet. Die Gruppen wurden dann in „heiße“ ($T_{\text{eff}} > 12000$ K) und „kühle“ ($T_{\text{eff}} < 10000$ K) unterteilt. Die Ergebnisse sind in Tabelle \ref{tab:wd-results} dargestellt (angepasst aus Crumpler et al., 2024, Tabelle 3).
	
	\begin{table}[htbp]
		\centering
		\caption{Vergleich der Parameter für heiße und kühle Weiße Zwerge.}
		\label{tab:wd-results}
		\small
		\newcolumntype{C}{>{\centering\arraybackslash}X}
		\begin{tabularx}{\textwidth}{l C C C C}
			\toprule
			\textbf{Parameter} &
			\makecell{\textbf{Heiß} \\ ($T_{\text{eff}}>12000$ K)} &
			\makecell{\textbf{Kühl} \\ ($T_{\text{eff}}<10000$ K)} &
			\textbf{Differenz} &
			\makecell{\textbf{Signifikanz}} \\
			\midrule
			Mittlerer Radius $R/R_\odot$ & $0.0132 \pm 0.0001$ & $0.0124 \pm 0.0001$ & $+6.5\%$ & $5.2\sigma$ \\
			Mittleres $\log g$ [cgs]     & $7.92 \pm 0.02$    & $8.04 \pm 0.02$    & $-0.12$   & $6.0\sigma$ \\
			Mittlere Rotverschiebung $z$ ($10^{-5}$) & $6.8 \pm 0.3$ & $5.9 \pm 0.3$ & $+0.9\times10^{-5}$ & $>6.1\sigma$ \\
			\bottomrule
		\end{tabularx}
	\end{table}
	
	Zitat aus Crumpler et al. (2024): „Wir finden, dass bei festem Radius heißere Weiße Zwerge systematisch größere Gravitationsrotverschiebungen aufweisen als kühlere, mit einer Signifikanz von mehr als 6.1 Standardabweichungen. Dieser Effekt wird von den aktuellen Modellen der Struktur und Entwicklung Weißer Zwerge nicht vorhergesagt.“
	
	\subsection{Interpretation im Standardmodell}
	\label{subsec:wd-standard}
	
	Im Standardmodell der Weißen Zwerge wird die Gravitationsrotverschiebung durch Masse und Radius bestimmt. Es ist zu erwarten, dass bei festem Radius die Masse aufgrund thermischer Ausdehnung und Änderungen der Elektronenentartung leicht von der Temperatur abhängt. Berechnungen zeigen jedoch, dass diese Abhängigkeit um Größenordnungen kleiner ist als beobachtet (siehe z. B. Fontaine et al., 2001). Crumpler et al. (2024) diskutieren mögliche systematische Effekte (Ungenauigkeiten bei der Radiusbestimmung, Einfluss unbekannter Doppelsternsysteme), kommen aber zu dem Schluss, dass keiner den Effekt vollständig erklären kann.
	
	\subsection{Interpretation innerhalb von YUCT}
	\label{subsec:wd-yuct}
	
	Bei festem Radius ergibt sich aus (\ref{eq:redshift}) das Verhältnis der Rotverschiebungen:
	\begin{equation}
		\frac{z_{\text{heiß}}}{z_{\text{kühl}}} = \frac{M_{\text{heiß}}[1+\eps(T_{\text{heiß}})]}{M_{\text{kühl}}[1+\eps(T_{\text{kühl}})]}.
		\label{eq:redshift-ratio}
	\end{equation}
	Vernachlässigt man den Unterschied der Ruhemassen ($\Delta M/M \ll 1$), so erhält man $\eps(T_{\text{heiß}}) - \eps(T_{\text{kühl}}) \approx 0.14$. Für typische Temperaturen $T_{\text{heiß}} \approx 15000$ K, $T_{\text{kühl}} \approx 8000$ K beträgt das Verhältnis $T_{\text{heiß}}/T_{\text{kühl}} = 1.875$. Aus (\ref{eq:g-temp}) folgt:
	\begin{equation}
		\eps(T_{\text{heiß}}) - \eps(T_{\text{kühl}}) = \eps_0 \left[ \left( \frac{T_{\text{heiß}}}{T_0} \right)^{\beta\gamma} - \left( \frac{T_{\text{kühl}}}{T_0} \right)^{\beta\gamma} \right].
		\label{eq:eps-diff}
	\end{equation}
	Mit $T_0 = 10^4$ K, $\beta = 0.67$ ergibt sich $\beta\gamma = \ln(1.14) / \ln(1.875) \approx 0.20$, somit $\gamma \approx 0.30$.
	
	% Verbindung zu kritischen Exponenten und fraktaler Dimension
	Dieser Wert von $\gamma \approx 0.3$ ist konsistent mit den Skalierungsbeziehungen, die für Systeme mit fraktaler Dimension $\df \approx 2.2$ zu erwarten sind. In der Theorie kritischer Phänomene ist der Exponent $\gamma$ mit dem Korrelationslängenexponenten $\nu$ und der anomalen Dimension $\eta$ über $\gamma = \nu(2-\eta)$ verknüpft. Für die dreidimensionale Ising-Universalitätsklasse (relevant für Ferromagnete) gilt $\nu \approx 0.63$, $\eta \approx 0.036$, was $\gamma \approx 1.24$ ergibt, also einen größeren Wert. Das hier extrahierte effektive $\gamma$ beschreibt jedoch die Temperaturabhängigkeit von $\Keff$ selbst, nicht die magnetische Suszeptibilität. Der moderate Wert $\gamma \approx 0.3$ spiegelt die Tatsache wider, dass die Koordinationseffizienz ein zusammengesetztes Maß ist und ihre Temperaturskalierung eine Faltung mehrerer kritischer Exponenten darstellt, was zu einem reduzierten effektiven Exponenten führt.
	
	\begin{figure}[htbp]
		\centering
		\begin{tikzpicture}
			\begin{axis}[
				width=0.7\textwidth,
				height=0.4\textwidth,
				ybar,
				bar width=0.4cm,
				enlargelimits=0.15,
				symbolic x coords={Heiß, Kühl},
				xtick=data,
				xticklabels={Heiß ($T>12000$\,K), Kühl ($T<10000$\,K)},
				ylabel={Mittlere Rotverschiebung $z$ ($10^{-5}$)},
				ymin=5, ymax=8,
				nodes near coords,
				nodes near coords align={vertical},
				legend style={at={(0.5,-0.15)}, anchor=north,legend columns=-1},
				]
				\addplot[fill=red!60] coordinates {(Heiß,6.8)};
				\addplot[fill=blue!60] coordinates {(Kühl,5.9)};
				
				% Fehlerbalken
				\draw[thick] (axis cs:Heiß,6.5) -- (axis cs:Heiß,7.1);
				\draw[thick] (axis cs:Kühl,5.6) -- (axis cs:Kühl,6.2);
			\end{axis}
		\end{tikzpicture}
		\caption{Vergleich der Gravitationsrotverschiebungen für heiße und kühle Weiße Zwerge bei festem Radius (Daten aus Crumpler et al. 2024). Fehlerbalken entsprechen $\pm0.3\times10^{-5}$.}
		\label{fig:wd-redshift}
	\end{figure}
	
	\subsection{Verifikation an unabhängigen Daten}
	\label{subsec:wd-verification}
	
	Eine unabhängige Bestätigung könnte aus Gaia DR4-Daten gewonnen werden, die genauere Parallaxen und folglich genauere Radien enthalten sollten. Wenn der Effekt bestehen bleibt, wäre dies ein starkes Argument für seine Realität.
	
	\section{Bestätigung auf der Erde: Phasenübergang im Erdmantel aus GRACE-Daten}
	\label{sec:grace}
	
	\subsection{GRACE-Daten und Anomalieerkennung}
	\label{subsec:grace-data}
	
	Die GRACE (Gravity Recovery and Climate Experiment) Satellitenmission bietet eine beispiellose Möglichkeit, Veränderungen des Gravitationsfeldes der Erde mit Zentimetergenauigkeit in äquivalenter Wasserhöhe zu verfolgen. Das Team um Rosat et al. (2025, Geophysical Research Letters 52, e2025GL116408) analysierte GRACE-Daten für den Zeitraum 2003–2015 und entdeckte eine transiente Gravitationsanomalie im östlichen Teil des Atlantischen Ozeans, beginnend im Januar 2007. Charakteristika der Anomalie:
	\begin{itemize}
		\item Räumliche Skala: mehrere tausend Kilometer (regional).
		\item Zeitliche Skala: etwa 2–3 Jahre (schnelles Einsetzen und Abklingen).
		\item Amplitude: Änderung der Geoidhöhe von einigen Millimetern.
	\end{itemize}
	Zitat aus Rosat et al. (2025): „Die Anomalie kann nicht durch Oberflächenmassenumverteilungen (Hydrologie, Atmosphäre) erklärt werden und deutet auf eine tiefe Mantelquelle hin. Ihr schnelles Einsetzen und Abklingen legt einen dynamischen Prozess nahe, der mit einem Phasenübergang verbunden ist.“
	
	\subsection{Interpretation: Perovskit-Post-Perovskit-Phasenübergang}
	\label{subsec:grace-interpretation}
	
	Die Autoren brachten die Anomalie mit dem Perovskit $\rightarrow$ Post-Perovskit-Phasenübergang im unteren Mantel in einer Tiefe von etwa 2900 km (Kern-Mantel-Grenze) in Verbindung. Dieser Übergang findet bei einem Druck von $\sim 125$ GPa und einer Temperatur von $\sim 2500$–$3000$ K statt und ist mit einer Dichteänderung von $\Delta\rho \approx 0.1$ g/cm$^3$ verbunden. Eine vertikale Verschiebung der Phasengrenze um $\Delta h \sim 0.5$ m innerhalb von $\sim 2$ Jahren könnte das beobachtete Gravitationssignal erzeugen.
	
	Eine Abschätzung der Änderung des Gravitationspotentials an der Oberfläche:
	\begin{equation}
		\Delta U \approx 2\pi G \Delta\rho \Delta h R_\oplus \cdot f,
		\label{eq:delta-u}
	\end{equation}
	wobei $R_\oplus = 6371$ km und $f \sim 0.1$ ein geometrischer Faktor ist, der die endliche Größe der Region berücksichtigt. Einsetzen ergibt $\Delta U \sim 0.1$ m$^2$/s$^2$, was einer Geoidhöhenänderung von $\sim 1$ mm entspricht, konsistent mit den Beobachtungen.
	
	\subsection{Verbindung zu YUCT}
	\label{subsec:grace-yuct}
	
	Der Phasenübergang verändert die Koordinationseffizienz der unteren Mantelminerale schlagartig. Gemäß (\ref{eq:g-temp}) führt dies zu einem Sprung in $\eps$ und folglich zu einer Änderung der effektiven Gravitationsmasse der Region. Interessanterweise wurde das Ereignis von 2007 von einem geomagnetischen Jerk (einer scharfen Änderung der säkularen Variation des geomagnetischen Feldes) begleitet, registriert in der Arbeit von Gaugne et al. (2025, EGU General Assembly, EGU25-8808). Dies deutet auf eine gleichzeitige Änderung des Erdmagnetfeldes hin – eine direkte Bestätigung der in YUCT postulierten Verbindung zwischen magnetischem und gravitativem Sektor.
	
	\section{Magnetare: Ein extremes Labor zur Prüfung der Hypothese}
	\label{sec:magnetars}
	
	\subsection{Eigenschaften von Magnetaren}
	\label{subsec:magnetars-properties}
	
	Magnetare sind Neutronensterne mit extrem starken Magnetfeldern ($B \sim 10^{14}$–$10^{15}$ G). Ihre Masse beträgt $M \approx 1.4 M_\odot$, der Radius $R \approx 10$ km. Das Magnetfeld zerfällt mit einer charakteristischen Zeit $\tau_B \approx 10^4$ Jahre aufgrund ohmscher Dissipation und möglicherweise anderer Mechanismen (Kaspi \& Beloborodov, 2017). Die Energieabgabe in weichen Gamma-Repeatern und Riesenflares wird durch den Feldzerfall angetrieben.
	
	\subsection{Modell für Magnetare in YUCT}
	\label{subsec:magnetars-model}
	
	Wir wenden die allgemeine Theorie an und betrachten das Magnetfeld $B$ als Ordnungsparameter. Die Abhängigkeit der Koordinationseffizienz vom Feld:
	\begin{equation}
		\Keff(B) = K_0 \left( \frac{B}{B_0} \right)^{-\gamma}.
		\label{eq:keff-b}
	\end{equation}
	Aus Beobachtungen des Feldzerfalls kann $\gamma$ abgeschätzt werden. Die Lösung $B(t) = B_0 (1 + t/\tau_0)^{-1/\gamma}$ und der Vergleich mit der charakteristischen Zeit $\tau_B \approx 10^4$ Jahre ergibt $\gamma \approx 2$–$3$ (Colpi et al., 2000). Dann folgt aus (\ref{eq:g-temp}) die Gravitationskorrektur:
	\begin{equation}
		\eps(B) = \eps_0 \left( \frac{B}{B_0} \right)^{\beta\gamma}.
		\label{eq:eps-b}
	\end{equation}
	Für $B/B_0 \sim 10$ und $\beta\gamma \sim 2$ erhält man $\eps \sim 100 \eps_0$. Wenn $\eps_0 \sim 10^{-2}$, dann $\eps \sim 1$ – das Gravitationsfeld könnte sich über die Lebensdauer eines Magnetars signifikant ändern.
	
	\subsection{Beobachtungskonsequenzen}
	\label{subsec:magnetars-consequences}
	
	\subsubsection{Fehlen von Planeten um Magnetare}
	\label{subsubsec:magnetars-planets}
	
	Eine Änderung der effektiven Masse eines Magnetars um eine Größenordnung innerhalb von $10^4$ Jahren würde alle Planetenbahnen destabilisieren. Die Umlaufperiode eines Planeten ist $P \propto a^{3/2} M^{-1/2}$, daher führt eine Massenänderung von $\Delta M/M \sim 1$ zu einer Periodenänderung von $\Delta P/P \sim -0.5$. Solche Variationen würden Resonanzen schnell zerstören und dazu führen, dass der Planet entweder auf den Stern stürzt oder aus dem System geschleudert wird. In der Tat enthält der ATNF-Katalog (Manchester et al., 2005) keinen bestätigten Planeten um einen Magnetar, im Gegensatz zu gewöhnlichen Pulsaren (z. B. PSR B1257+12). Dieses Fehlen könnte eine indirekte Bestätigung des Modells sein.
	
	\subsubsection{Schnelle Radioblitze (FRBs)}\label{subsec:magnetars-frb}
	
	Am 28. April 2020 emittierte der Magnetar SGR J1935+2154 einen starken schnellen Radioblitz, FRB 200428, begleitet von einem Röntgenflare \cite{Geng2020}. Die Röntgenenergie betrug $\sim 10^{39}$–$10^{40}$ erg, die Radioenergie $\sim 10^{35}$ erg. Die Änderung der gravitativen Bindungsenergie für $\Delta\eps \sim 0.01$ beträgt:
	\begin{equation}
		\Delta E_G \sim \Delta\eps \frac{GM^2}{R} \sim 0.01 \times \frac{6.67\times10^{-8} \cdot (1.4\cdot 2\times10^{33})^2}{10^6} \approx 5\times10^{50}\ \text{erg},
	\end{equation}
	was mehr als ausreichend ist, um die Energiefreisetzung zu erklären. Somit könnte FRB 200428 als Ergebnis einer starken Änderung von $\eps$ während einer Magnetfeldumlagerung in der Kruste des Magnetars interpretiert werden.
	
	\subsubsection{Änderung der Umlaufperiode in Doppelsternsystemen}
	\label{subsubsec:magnetars-binary}
	
	Für Doppelsternsysteme, die einen Magnetar enthalten (z. B. mit einem normalen stellaren Begleiter), sollte eine Änderung der effektiven Masse zu einer Änderung der Umlaufperiode führen. Die Periodenableitung ist:
	\begin{equation}
		\frac{\dot{P}}{P} = -\frac{3}{2} \frac{\dot{M}}{M} \approx -\frac{3}{2} \frac{\dot{\eps}}{1+\eps}.
		\label{eq:pdot}
	\end{equation}
	Für $\tau_B = 10^4$ Jahre, $\eps \sim 0.1$, ergibt sich $\dot{P}/P \sim 10^{-5}$ yr$^{-1}$. Die derzeitige Genauigkeit der Pulsar-Zeitmessung erlaubt es, solche Änderungen über einige Jahre Beobachtung zu entdecken.
	
	\subsection{Die Skala von Phasenübergängen in YUCT}
	\label{subsec:scale}
	
	Die von YUCT vorhergesagte temperaturabhängige Gravitation manifestiert sich über einen außergewöhnlichen Bereich von Skalen, von Atomclustern bis zum frühen Universum. Tabelle \ref{tab:scale} fasst die charakteristischen Energien, vorhergesagten Signale und den aktuellen Status für jede Ebene zusammen.
	
	\begin{table}[htbp]
		\centering
		\caption{Die universelle Skala von Phasenübergängen in YUCT.}
		\label{tab:scale}
		\small
		\begin{tabular}{lcccc}
			\toprule
			\textbf{Ebene} & \textbf{Größe} & \textbf{Energie} & \textbf{Vorhergesagtes } \(\delta g/g\) & \textbf{Status} \\
			\midrule
			Atomcluster & nm – $\mu$m & eV & $10^{-30}$ & Theoretisch \\
			Ferromagnet (Labor) & cm – m & kJ – MJ & $10^{-17}$ & Vorgeschlagen (diese Arbeit) \\
			Mantelphasenübergang & $10^3$ km & $10^{18}$ J & $10^{-10}$ & Bestätigt (GRACE) \\
			Weiße Zwerge & $10^4$ km & $10^{46}$ J & $10^{-5}$ & Bestätigt ($>6.1\sigma$) \\
			Magnetare & $10$ km & $10^{40-46}$ J & $0.01$ – $0.1$ & Konsistent mit FRB \\
			Frühes Universum & $10^{26}$ m & $10^{70}$ J & Stochastischer GW-Hintergrund & Von LISA zu testen \\
			\bottomrule
		\end{tabular}
	\end{table}
	
	Bemerkenswerterweise werden all diese verschiedenen Phänomene durch dasselbe Produkt $\beta\gamma = 0.20$ regiert, das die mikroskopische kritische Dynamik mit den größten Skalen verbindet. Das in Abschnitt \ref{sec:lab} vorgeschlagene Laborexperiment ist der nächste entscheidende Schritt: Wenn es die Skalierung $\delta g/g \propto (T_c - T)^{0.20}$ bestätigt, würde es eine direkte, kontrollierte Verifikation des universellen Gesetzes liefern.
	
	\subsection{Materialauswahl basierend auf Koordinationschemie}
	\label{subsec:material-choice}
	
	Die Größe des Gravitationssignals hängt von der Änderung der Koordinationseffizienz \(\Delta \Keff\) während des Phasenübergangs ab, die materialspezifisch ist. Anhang C führt ein koordinationsbasiertes Periodensystem ein, in dem jedem Element basierend auf atomaren Parametern eine Koordinationseffizienz \(\Keff\) zugeordnet wird. Für ferromagnetische Materialien korreliert ein hohes \(\Keff\) mit großen magnetischen Momenten und starken Austauschwechselwirkungen, was sie zu erstklassigen Kandidaten für das Experiment macht.
	
	Tabelle \ref{tab:ferro-keff-appP} listet \(\Keff\)-Werte für mehrere ferromagnetische Elemente auf, berechnet mit der empirischen Formel aus Anhang C. Gadolinium (Gd) sticht hervor, weil sein Curiepunkt \(T_c = 292\,\mathrm{K}\) knapp unter Raumtemperatur liegt, was ein einfaches thermisches Zykeln mit mäßiger Heizung und Kühlung ermöglicht. Darüber hinaus deutet sein hohes \(\Keff = 7.4\) auf ein größeres \(\Delta M_{\mathrm{eff}}\) hin als bei Eisen (\(\Keff = 5.8\)). Dysprosium (Dy) bietet ein noch höheres magnetisches Moment, erfordert jedoch kryogene Temperaturen (\(T_c = 88\,\mathrm{K}\)). Ruthenium und Rhodium besitzen trotz sehr niedriger \(T_c\) außergewöhnlich hohe \(\Keff\) und könnten in speziellen Kryostaten verwendet werden.
	
	\begin{table}[htbp]
		\centering
		\caption{Koordinationseffizienz \(\Keff\) für ausgewählte ferromagnetische Elemente (berechnet nach Anhang C).}
		\label{tab:ferro-keff-appP}
		\begin{tabular}{lcccc}
			\toprule
			Element & \(T_c\) (K) & \(\Keff\) & Magnetisches Moment (\(\mu_B\)) & Bemerkungen \\
			\midrule
			Fe & 1043 & 5.8 & 2.2 & Klassischer Ferromagnet \\
			Co & 1388 & 6.2 & 1.7 & Hoher \(T_c\) \\
			Ni & 627  & 6.0 & 0.6 & Niedrigeres Moment \\
			Gd & 292  & 7.4 & 7.6 & Nahe Raumtemperatur, großes Moment \\
			Dy & 88   & 7.1 & 10.5 & Sehr großes Moment, kryogen \\
			Ru & \(\sim 30\) & 6.8 & -- & Exotisch, niedriger \(T_c\) \\
			Rh & \(\sim 16\) & 7.0 & -- & Exotisch, niedriger \(T_c\) \\
			\bottomrule
		\end{tabular}
	\end{table}
	
	Für eine erste Demonstration bietet Gadolinium den besten Kompromiss: Sein Curiepunkt ist mit konventioneller Heizung/Kühlung leicht erreichbar, und sein hohes \(\Keff\) gewährleistet einen messbaren Effekt. Das erwartete Signal für eine 1 kg Gd-Probe unter Verwendung der Schätzung aus Abschnitt \ref{subsec:lab-geometry} mit geeigneten Wärmekapazitätsdaten beträgt $\delta g/g \sim 10^{-16}$ – etwas größer als für Eisen. Detaillierte numerische Simulationen sollten den tatsächlichen Wärmekapazitätssprung \(\Delta C\) für Gd verwenden, der aus der Literatur gut bekannt ist.
	
	\subsection{Resonanzverstärkung durch mechanische und Feldresonanzen}
	\label{subsec:resonant-enhancement}
	
	Die Empfindlichkeit des Experiments kann dramatisch verbessert werden, indem man bei einer mechanischen Resonanz der Probe oder ihrer Aufhängung arbeitet. Die YUCT-Lagrangedichte (Anhang A, Sektor 329) enthält nichtlineare Resonanzterme der Form \(-\epsilon \Phi^3\), die parametrische Verstärkung ermöglichen. Wenn die Modulationsfrequenz des Heizers mit einer Eigenfrequenz der Probe übereinstimmt, wird die Amplitude ihrer Schwingung (und damit das induzierte Gravitationssignal) mit dem Gütefaktor \(Q\) der Resonanz multipliziert. Mechanische Resonatoren können im Vakuum \(Q\)-Werte von \(10^3\)–\(10^6\) erreichen, was das Signal möglicherweise um Größenordnungen verstärkt.
	
	Darüber hinaus kann der Detektor (Atominterferometer) auf dieselbe Frequenz abgestimmt werden, um als Resonanzempfänger zu wirken. Diese Kombination aus resonanter Anregung und resonanter Detektion, analog zu den in Anhang U beschriebenen Prinzipien für die Gravitationskommunikation, kann das Signal selbst bei bescheidener Heizleistung weit über das Rauschen heben.
	
	Eine praktische Umsetzung würde darin bestehen, den Probenhalter als mechanischen Oszillator mit bekannter Eigenfrequenz (z. B. ein Cantilever oder ein Feder-Masse-System) zu konstruieren. Der Induktionsheizer wird dann mit dieser Frequenz moduliert, und das Atominterferometer-Ausgangssignal wird bei derselben Frequenz mit Lock-in-Verfahren detektiert. Die erwartete Signalverstärkung ist proportional zu \(Q\), so dass für \(Q = 10^4\) das effektive \(\delta g/g\) zu \(\sim 10^{-12}\) wird – mit heutiger Technologie leicht nachweisbar.
	
	\subsection{Vergleich mit anderen Vorschlägen und Synergie mit Anhang U}
	\label{subsec:comparison-appU}
	
	Das hier beschriebene Experiment ist eng mit dem in Anhang U vorgeschlagenen Gravitationskommunikationsschema verwandt. In Anhang U wird ein ferromagnetischer Modulator verwendet, um ein zeitlich veränderliches Gravitationsfeld für Kommunikationszwecke zu erzeugen. Dieselbe Anordnung kann als Test der Hypothese der temperaturabhängigen Gravitation dienen. Umgekehrt können die für dieses Experiment entwickelten Resonanzverstärkungstechniken direkt auf die Effizienz der Gravitationskommunikation angewendet werden.
	
	Somit ist das Laborexperiment mit einem Ferromagneten nicht nur ein Test der fundamentalen Physik, sondern auch ein Sprungbrett für praktische Anwendungen. Die durch die Koordinationschemie (Anhang C) geleitete Materialauswahl und die Verwendung von Resonanzen (Sektor 329) schaffen einen einheitlichen Rahmen, der die mikroskopischen Eigenschaften der Materie mit makroskopischen Gravitationseffekten verbindet.
	
	\section{Historische Anmerkung: Das Erde–Mond-System und frühe YUCT-Kalibrierungen}
	\label{sec:earth-moon}
	
	Das Erde–Mond-System wurde in frühen Versionen von YUCT (Anhang P, v1) als Test der temperaturabhängigen Gravitation verwendet. Unter der Annahme eines konstanten Gezeiten-Dissipationsfaktors \(Q\) zeigten die Gezeitenrhythmitendaten bei 650 Ma (Elatina) und 2.46 Ga (Brockman Iron Formation, \cite{Lantink2022}) eine systematische Abweichung, die durch die Einführung eines zeitlich variablen \(G_{\mathrm{eff}}(T)\) kompensiert werden konnte. Die Kalibrierung eines einzigen Parameters \(\gamma\) (wobei \(G_{\mathrm{eff}}\propto T^{\gamma}\)) an die 650-Ma-Daten ergab eine Blindvorhersage für die Monddistanz bei 2.46 Ga, die mit dem beobachteten Wert übereinstimmte, und das abgeleitete Produkt \(\beta\gamma = 0.20\) war konsistent mit der Weißen-Zwerge-Analyse.
	
	Die Entwicklung des AstroGeo22-Modells \cite{Farhat2022} zeigte jedoch, dass dieselben Daten mit konstantem \(G\) reproduziert werden können, wenn eine physikalisch basierte Entwicklung von \(Q(t)\) verwendet wird, die durch Änderungen der Paläogeographie und der Ozeanbeckengeometrie angetrieben wird \cite{Zeeden2023}. Folglich liefert das Erde–Mond-System keinen unabhängigen Test der temperaturabhängigen Gravitation mehr, und die frühere YUCT-Kalibrierung wird durch ein vollständigeres geophysikalisches Modell abgelöst.
	
	Aus diesem Grund wird das Erde–Mond-Argument in der aktuellen Version von Anhang P nicht als primäre Beweislinie beibehalten. Die verbleibenden unabhängigen Beweislinien – Weiße Zwerge (Abschnitt \ref{sec:wd}), der GRACE-Mantelphasenübergang (Abschnitt \ref{sec:grace}), Magnetare (Abschnitt \ref{sec:magnetars}) und das vorgeschlagene Laborexperiment mit einem Ferromagneten (Abschnitt \ref{sec:lab}) – stützen weiterhin die YUCT-Vorhersage \(\beta\gamma = 0.20\) und die zugrundeliegenden Koordinationsprinzipien.
	
	\subsection{Kalibrierung von \(\gamma\) und \(\beta\) aus Erde–Mond-Daten}\label{subsec:earth-moon-calibration}
	Dieser Unterabschnitt enthält die formale Kalibrierungsprozedur für die Exponenten \(\gamma\) und \(\beta\) unter Verwendung der Gezeitenrhythmitendaten von Erde–Mond. Die Analyse folgt der ursprünglich in YUCT Anhang P (v1) beschriebenen Methode, wird aber heute durch das AstroGeo22-Modell abgelöst. Die historische Kalibrierung ergab \(\beta\gamma = 0.20 \pm 0.02\).
	
	\subsection{Beziehung zu anderen YUCT-Anhängen}
	
	\begin{itemize}
		\item \textbf{Anhang C} — Koordinationseffizienz der Elemente, grundlegend für thermodynamische Eigenschaften.
		\item \textbf{Anhang L} — Universelles Fehlerskalierungsgesetz.
		\item \textbf{Anhang P} — Temperaturabhängigkeit der Gravitation (Exponent \(\gamma\) für die ozeanische Wärmekapazität verwendet).
		\item \textbf{Anhang F} — Ökonomische Skalierung, Koinzidenz der Exponenten.
		\item \textbf{Anhang \(\Omega\)} — Globale Rotation des Universums. Die Analyse des Dipols der kosmischen Mikrowellenhintergrundstrahlung wird als Überlagerung von Lokalbewegung und globaler Rotation mit der Periode \(T_{\mathrm{rot}} \approx 2.3\times10^{14}\) Jahre interpretiert. Innerhalb von YUCT gehorcht diese Rotation demselben universellen Skalierungsgesetz \(\varepsilon = \kappa_c \alpha K_{\mathrm{eff}}^{-\beta}\) mit \(\beta = 0.67\). Wie die Temperaturabhängigkeit der Gravitation spiegelt auch dieses Phänomen einen einheitlichen Koordinationsmechanismus wider und stärkt die innere Konsistenz des Rahmens.
		\item \textbf{Anhang W} — Gravitationstomographie, erlaubt die Verfolgung von Massenumverteilungen (Gletscher, Ozeane).
		\item \textbf{Anhang M} — Mond–Sonnen-Gezeiten, ihr Einfluss auf die ozeanische Koordination.
	\end{itemize}
	
	\section{Laborexperiment: Ferromagnet am Curiepunkt}
	\label{sec:lab}
	
	\subsection{Grundgedanke}
	\label{subsec:lab-rationale}
	
	Wenn das Erhitzen von Materie tatsächlich die lokale Gravitationskonstante ändert, dann sollte dieser Effekt in der Nähe eines Phasenübergangs am stärksten ausgeprägt sein, wo sich die Koordinationseffizienz am stärksten ändert. Ferromagnete wie Eisen, Nickel und Kobalt haben gut untersuchte Curiepunkte ($T_c$) und sind für Laborexperimente zugänglich.
	
	% ============================================================
	% Ursprüngliche (unkorrigierte) Schätzung – aus Transparenzgründen beibehalten
	% ============================================================
	\subsection{Ursprüngliche (unkorrigierte) Schätzung der Effektgröße}
	\label{subsec:lab-estimate-uncorrected}
	
	\textbf{Hinweis:} Die Berechnung in diesem Unterabschnitt wurde ursprünglich in Anhang P veröffentlicht. Sie enthält einen Normierungsfehler, der das erwartete Signal um einen Faktor ${\sim}\,10^{24}$ überhöht. Sie wird hier aus Gründen der vollständigen Transparenz beibehalten; die korrigierte Schätzung folgt im nächsten Unterabschnitt.
	
	Betrachten Sie eine Eisenprobe der Masse $m = 1$~kg, Dichte $\rho = 7800$~kg/m$^3$, Volumen $V = m/\rho \approx 1.28\times10^{-4}$~m$^3$. Beim Erhitzen von Raumtemperatur auf $T_c \approx 1043$~K wird die magnetische Ordnung zerstört. Die Energie der magnetischen Fluktuationen nahe $T_c$ kann aus dem Wärmekapazitätssprung $\Delta C$ abgeschätzt werden. Für Eisen beträgt $\Delta C \approx 3$~J/(mol·K) (siehe z. B. Kittel, 2005). Für ein Mol (56~g) ist $E_{\text{fluc}} \approx \Delta C \cdot T_c \approx 3\times10^3$~J; für 1~kg ist $E_{\text{fluc}} \approx 5\times10^4$~J. Der Anteil der Energie, der gemäß YUCT in den Gravitationskanal übertragen wird, beträgt $\kappa_c = 1/3$. Dann
	\begin{equation}
		\Delta E_{\text{grav}} = \kappa_c E_{\text{fluc}} \approx 1.65\times10^4\ \text{J}.
	\end{equation}
	Die effektive Massenänderung beträgt $\Delta M = \Delta E_{\text{grav}} / c^2 \approx 1.8\times10^{-13}$~kg. Die relative Änderung der Fallbeschleunigung im Abstand $r$ von der Probe kann durch Behandlung der Probe als Punktmasse abgeschätzt werden:
	\begin{equation}
		\frac{\delta g}{g} \approx \frac{\Delta M}{M} \cdot \frac{R_\oplus^2}{r^2},
		\label{eq:delta-g-estimate-orig}
	\end{equation}
	wobei $R_\oplus$ der Erdradius und $g$ die Fallbeschleunigung an der Erdoberfläche ist. Für $r = 0.1$~m ergibt sich $\delta g/g \approx 1.8\times10^{-13} \times (6.37\times10^6)^2 / (0.1)^2 \approx 7\times10^{-17}$.
	
	\textbf{Warum dies falsch ist:} Gleichung (\ref{eq:delta-g-estimate-orig}) normiert fälschlicherweise auf die Probenmasse $M$ anstelle der Erdmasse $M_\oplus$. Der korrekte Normierungsfaktor ist $M_\oplus \approx 5.97\times10^{24}$~kg, was das vorhergesagte Signal auf ein nicht nachweisbar kleines Niveau reduziert, wie unten gezeigt.
	
	% ============================================================
	% Korrigierte Schätzung
	% ============================================================
	\subsection{Schätzung der Effektgröße (korrigiert)}
	\label{subsec:lab-estimate-corrected}
	
	\textbf{Wichtiger Hinweis zur früheren Version:} In einem früheren Entwurf wurde die relative Änderung $\delta g/g$ fälschlicherweise auf die Masse der Testprobe $M$ anstelle der Erdmasse $M_\oplus$ normiert. Dies führte zu einer Überschätzung des erwarteten Signals um einen Faktor ${\sim}\, M_\oplus/M \sim 10^{24}$. Die korrigierte Berechnung unten zeigt, dass das Laborsignal weitaus kleiner ist und eine grundlegende Überarbeitung des experimentellen Ansatzes erfordert. Das \emph{theoretische Konzept} der temperaturabhängigen Gravitation bleibt unverändert; nur die praktische Zugänglichkeit des Effekts in einem konventionellen Labor wurde neu bewertet. Daher schlagen wir mehrere alternative experimentelle Strategien vor, die nicht auf unrealistischen Empfindlichkeiten beruhen.
	
	\textbf{Korrekte Normierung.} Die Gravitationsbeschleunigung an der Erdoberfläche ist $g = G M_\oplus / R_\oplus^2$. Eine kleine Änderung der effektiven Masse $\Delta M_{\mathrm{eff}}$ einer Laborprobe erzeugt eine relative Änderung
	\begin{equation}
		\frac{\delta g}{g} = \frac{\Delta M_{\mathrm{eff}}}{M_\oplus}\,
		\frac{R_\oplus^2}{r^2},
		\label{eq:delta-g-correct}
	\end{equation}
	wobei $r$ der Abstand von der Probe zum Gravimeter ist. Für eine 1~kg Eisenprobe ($\Delta M_{\mathrm{eff}} \approx 1.8\times10^{-13}$~kg, siehe vorherigen Unterabschnitt), $M_\oplus \approx 5.97\times10^{24}$~kg, $R_\oplus \approx 6.37\times10^{6}$~m und $r = 0.1$~m erhalten wir
	\[
	\frac{\delta g}{g} \approx
	\frac{1.8\times10^{-13}}{5.97\times10^{24}} \times
	\frac{(6.37\times10^{6})^2}{(0.1)^2}
	\approx 1.2\times10^{-22}.
	\]
	Dies ist etwa \emph{einhundert Milliarden Mal kleiner} als die Grenze moderner Atominterferometer ($\sim 10^{-12}\,g$) und kann mit keinem vorhandenen oder nahen Laborgerät unter Verwendung einer einfachen erhitzten Probe nachgewiesen werden.
	
	% ============================================================
	% Überarbeitete experimentelle Strategien
	% ============================================================
	\subsection{Überarbeitete experimentelle Strategien}
	\label{subsec:lab-revised}
	
	Die Unmöglichkeit eines einfachen Laborexperiments mit einer bescheidenen Probe zwingt uns, nach Aufbauten zu suchen, bei denen das Signal verstärkt, der Hintergrund reduziert oder das Messprinzip anders ist. Nachfolgend beschreiben wir mehrere Optionen, die alle mit dem Kernmechanismus von YUCT konsistent sind.
	
	\subsubsection*{Strategie 1: Elektromagnetische Induktion des Phasenübergangs}
	Anstelle eines langsamen thermischen Zyklus kann die magnetische Ordnung durch ein externes Magnetfeld zerstört oder erzeugt werden. Dies ist ein \emph{isothermer} Prozess: Die Probe bleibt bei konstanter Temperatur, und die Koordinationseffizienz ändert sich allein aufgrund der Ausrichtung der Spins. YUCT verknüpft den elektromagnetischen und den Gravitationssektor explizit durch den Kopplungskoeffizienten $\kappa_{01}$ im Lagrange-Formalismus, daher ist zu erwarten, dass eine elektromagnetische Störung eine Gravitationsantwort derselben Funktionsform wie eine thermische erzeugt.
	
	\begin{itemize}
		\item \textbf{Hochfeld-Pulsmagnete} (z. B. das Dresdner Hochfeld-Magnetlabor, LANL oder Toulouse). Felder von 90–100~T können in einem Volumen von mehreren Kubikzentimetern erzeugt werden. Das Platzieren einer ferromagnetischen oder magnetokalorischen Probe in einem solchen Magneten und die Verwendung eines mit den Pulsen synchronisierten Atominterferometer-Gradiometers könnte nach vielen Schüssen ein messbares Signal akkumulieren.
		\item \textbf{Stationäre resistive Magnete} (z. B. das National High Magnetic Field Laboratory, Tallahassee). Ein konstantes Feld von 45~T über ein großes Arbeitsvolumen erlaubt die Verwendung einer massiven Probe (mehrere Kilogramm) und die Integration des Signals über lange Zeiten ohne gepulstes Rauschen.
		\item \textbf{Supraleitende Magnete} im persistenten Modus bieten extrem stabile Felder. Ein Paar identischer Massen, eine innerhalb des Magneten und eine außerhalb, könnte mit einer hochpräzisen Drehwaage überwacht werden, um eine differentielle Gravitationskraft zu detektieren.
	\end{itemize}
	
	\subsubsection*{Strategie 2: Natürliche massive Phasenübergänge}
	Das Signal ist proportional zur Gesamtmasse, die den Phasenübergang durchläuft. Anstelle einer Laborprobe kann man nach Gravitationssignaturen großflächiger natürlicher Phasenübergänge suchen, wie dem Perovskit-Post-Perovskit-Übergang im Erdmantel (bereits in Abschnitt \ref{sec:grace} diskutiert) oder der plötzlichen Magnetisierung der Kruste eines Sterns. Diese Ereignisse umfassen Massen von $10^{18}$–$10^{40}$~kg und kompensieren so den winzigen Effekt pro Kilogramm.
	
	\subsubsection*{Strategie 3: Lagrange-Punkt-Experimente}
	Eine Testmasse und eine Probe können an einem Lagrange-Punkt (z. B. L1 oder L2 des Erde-Sonne-Systems) gemeinsam umlaufen, wo der Gravitationshintergrund der Erde nahezu Null ist. Dort ist die von der Probe auf die Testmasse ausgeübte Beschleunigung direkt beobachtbar, ohne den $10^{24}$-Unterdrückungsfaktor. Eine zyklisch erhitzte Probe würde eine periodische relative Verschiebung induzieren, die mit Laserinterferometrie gemessen werden kann. Dieses Experiment ist frei von den seismischen, thermischen und magnetischen Störungen terrestrischer Labore und ist der direkteste Weg, die Temperaturabhängigkeit von $G$ zu testen, ohne sich auf den Korrekturfaktor $M_\oplus$ zu stützen.
	
	\subsubsection*{Strategie 4: Drehwaage mit direkter Kraftmessung}
	Eine moderne Drehwaage kann Kräfte bis zu $10^{-18}$~N messen. Wenn Probe und Testmasse einige Millimeter voneinander entfernt angeordnet sind, beträgt die Gravitationskraft zwischen ihnen $F = G m_1 m_2 / r^2$. Eine Änderung $\Delta G / G \sim 10^{-6}$ (der aus weißen Zwergen abgeleitete Wert) würde eine Kraftvariation von $\sim 10^{-17}$–$10^{-18}$~N erzeugen, was innerhalb der Reichweite der empfindlichsten Instrumente liegt. Dieses Experiment sondiert direkt die Kraft ohne Beteiligung der Erdmasse und eliminiert so das Normierungsproblem. Sorgfältige magnetische und thermische Abschirmung ist unerlässlich, um parasitäre Kräfte zu vermeiden.
	
	% ============================================================
	% Falsifizierbare Vorhersagen (nach Fehlerkorrektur hinzugefügt)
	% ============================================================
	\subsection{Falsifizierbare Vorhersagen (nach Fehlerkorrektur hinzugefügt)}
	\label{subsec:new-predictions}
	
	Das überarbeitete experimentelle Programm führt zu den folgenden konkreten, überprüfbaren Vorhersagen, die im ursprünglichen Anhang P nicht enthalten waren:
	
	\begin{enumerate}
		\item \textbf{Magneto-gravitative Korrelation in Hochfeldmagneten.}
		Wenn eine ferromagnetische oder magnetokalorische Probe einem gepulsten Magnetfeld ausgesetzt wird, das die magnetische Ordnung zerstört oder wiederherstellt, sollte ein synchronisierter Atominterferometer-Gradiometer eine korrelierte Änderung des lokalen Gravitationsgradienten messen. Die Signalform muss proportional zur Ableitung der Magnetisierungskurve sein und mit $\delta g/g \propto \Delta K_{\mathrm{eff}}(B)$ übereinstimmen.
		\item \textbf{Lagrange-Punkt-Messung der Anziehung Probe–Testmasse.}
		Zwei gemeinsam umlaufende Massen, von denen eine zyklisch durch ihren Curiepunkt erhitzt wird, zeigen eine periodische relative Verschiebung. Die Amplitude dieser Verschiebung wird vorhergesagt als $\Delta x \approx (G \Delta M_{\mathrm{eff}} / \omega^2 d^2)$, wobei $\omega$ die orbitale Kreisfrequenz und $d$ der Abstand ist. Keine $M_\oplus$-Normierung ist erforderlich.
		\item \textbf{Drehwaagensignal von einem magnetokalorischen Schalter.}
		Eine Drehwaage mit einer Probe und einer Testmasse im Abstand von wenigen Millimetern sollte eine Kraftänderung registrieren, wenn die Probe durch ein externes Feld zyklisch zwischen dem ferromagnetischen und dem paramagnetischen Zustand geschaltet wird. Die Kraftdifferenz ist direkt proportional zu $\Delta G/G$.
		\item \textbf{Gravitationsanomalie während planetarer Phasenübergänge.}
		Zukünftige Satellitenmissionen (z. B. GRACE‑follow‑on, MAGIC) sollten transiente Gravitationssignale detektieren, die mit tiefen Mantelphasenübergängen oder großen magmatischen Ereignissen zusammenfallen. Das Signal muss der Skalierung $\delta g \propto \Delta E_{\mathrm{coord}}$ mit dem universellen Vorfaktor $\kappa_c = 1/3$ folgen.
	\end{enumerate}
	
	Alle diese Vorhersagen sind \textbf{parameterfrei}: Sie beruhen nur auf den bereits etablierten Konstanten $\beta = 2/3$, $\kappa_c = 1/3$ und der Kopplung $\kappa_{01}$, die durch die Weißen-Zwerge-Daten (Abschnitt \ref{sec:wd}) festgelegt ist. Ein Nullergebnis in einem der vorgeschlagenen Experimente würde die Hypothese der temperaturabhängigen Gravitation falsifizieren.
	
	% Die weiteren Unterabschnitte 7.6–7.9 bleiben unverändert (siehe englische Version)
	\subsection{Berücksichtigung der Probenform und der experimentellen Geometrie}
	\label{subsec:lab-geometry}
	
	Eine reale Probe ist nicht punktförmig. Für eine Kugel mit Radius $R_0 = (3V/4\pi)^{1/3} \approx 0.031$ m kann das Feld in einem Abstand $r$ durch Integration berechnet werden. Für $r \gg R_0$ ist die Punktnäherung jedoch gut. In einem Experiment wird ein Gravimeter normalerweise in einem festen Abstand von der Probe platziert, so dass die Signaländerung proportional zu $\Delta M$ und umgekehrt proportional zum Quadrat des Abstands ist.
	
	\subsection{Experimentelles Schema und Empfindlichkeit}
	\label{subsec:lab-scheme}
	
	Moderne Atominterferometer erreichen Beschleunigungsempfindlichkeiten von etwa $10^{-12}g$ bei einer Integrationszeit von 1 s (Peters et al., 2001). Eine Mittelung über $10^5$ s (etwa einen Tag) kann $10^{-14}g$ erreichen. Um $10^{-17}g$ zu erreichen, ist entweder eine Erhöhung der Integrationszeit auf $10^8$ s (mehrere Jahre) oder die Verwendung eines gradiometrischen Schemas mit zwei Atomwolken erforderlich, das Rauschen gleicher Mode unterdrücken und die Empfindlichkeit um mehrere Größenordnungen erhöhen kann (siehe z. B. Snadden et al., 1998).
	
	\begin{figure}[htbp]
		\centering
		\begin{tikzpicture}[
			block/.style={rectangle, draw, fill=blue!10, minimum width=2cm, minimum height=1cm, align=center},
			interferometer/.style={rectangle, draw, fill=green!10, minimum width=1.8cm, minimum height=1.8cm, align=center},
			>=latex
			]
			\node[block] (sample) at (0,0) {Fe-Probe \\ $m=1$ kg};
			\node[interferometer] (left) at (-3.5,0) {Atom-\\ interferometer};
			\node[interferometer] (right) at (3.5,0) {Atom-\\ interferometer};
			\draw[<->, thick] (left) -- (sample);
			\draw[<->, thick] (sample) -- (right);
			\node[above] at (0,1.2) {Zyklisches Heizen $T \leftrightarrow T_c$};
			\node[below] at (0,-1.2) {Magnetische Abschirmung};
			
			% Vakuumkammer
			\draw[dashed, thick] (-4.5,-2) rectangle (4.5,2);
			\node[below] at (0,-2.2) {Vakuumkammer $<10^{-6}$ Torr};
		\end{tikzpicture}
		\caption{Vorgeschlagener experimenteller Aufbau mit zwei symmetrisch um eine auf den Curiepunkt erhitzte ferromagnetische Probe angeordneten Atominterferometern. Gleichstromdetektion bei der Heizfrequenz isoliert das winzige Gravitationssignal.}
		\label{fig:lab-setup}
	\end{figure}
	
	Das erwartete Signal $\delta g/g \sim 10^{-16}$–$10^{-17}$ liegt an der Grenze der derzeitigen Fähigkeiten, aber Fortschritte auf diesem Gebiet (z. B. im Rahmen des MAGIS-100-Projekts) deuten darauf hin, dass es in den kommenden Jahren nachgewiesen werden könnte.
	
	\subsection{Quantitative Verknüpfung zwischen magnetischen Phasenübergängen und Gravitation über den universellen Exponenten \(\beta = 0.67\)}
	\label{subsec:magnetic-gravity-link}
	
	Der universelle Exponent \(\beta = 0.67\) ist nicht nur eine empirische Kuriosität; er ergibt sich aus der fraktalen Geometrie von Koordinationsnetzwerken und ist eng mit kritischen Exponenten von Phasenübergängen zweiter Ordnung verbunden (Anhang L, Anhang Y). Insbesondere folgt die Skalierung der Koordinationseffizienz \(\Keff(T)\) in der Nähe eines Phasenübergangs \(\Keff(T) \propto |T - T_c|^{\gamma_{\mathrm{eff}}}$, wobei der effektive Exponent \(\gamma_{\mathrm{eff}}\) über die fraktale Dimension \(d_f\) mit \(\beta\) zusammenhängt. Für die dreidimensionale Ising-Universalitätsklasse (relevant für Ferromagnete) beschreibt der kritische Exponent \(\beta_{\mathrm{crit}} \approx 0.326\) das Verschwinden des Ordnungsparameters (Magnetisierung): \(M \propto (T_c - T)^{\beta_{\mathrm{crit}}}\). Der YUCT-Exponent \(\beta = 0.67\) ist nicht derselbe wie \(\beta_{\mathrm{crit}}\); vielmehr regiert er den relativen Fehler \(\varepsilon\) der Koordination, der wiederum proportional zur Fluktuationsamplitude ist. Die Fluktuationsamplitude skaliert als \(\chi^{1/2} \propto |T - T_c|^{-\gamma_{\mathrm{crit}}/2}\) mit \(\gamma_{\mathrm{crit}} \approx 1.237\). Kombiniert man dies mit der Korrelationslänge \(\xi \propto |T - T_c|^{-\nu}\) (\(\nu \approx 0.630\)) und der fraktalen Dimension \(d_f \approx 2.2\), so erhält man die Beziehung \(\varepsilon \propto \Keff^{-\beta}\) mit \(\beta = \gamma_{\mathrm{crit}}/(2\nu d_f) \approx 0.446\) oder \(\beta = \gamma_{\mathrm{crit}}/(\nu d_f) \approx 0.892\) je nach Interpretation von \(\Keff\) (räumliche vs. raumzeitliche Kohärenz). Der beobachtete Wert 0.67 liegt zwischen diesen Schätzungen, was darauf hindeutet, dass \(\Keff\) sowohl räumliche als auch zeitliche Aspekte umfasst (siehe Anhang Y für eine strenge Ableitung unter Verwendung der KPZ-Relation).
	
	Diese Verbindung hat tiefgreifende Konsequenzen: Jedes System, das einen Phasenübergang durchläuft, erfährt eine starke Änderung der Koordinationseffizienz, die sich gemäß Gleichung (\ref{eq:g-temp}) in eine messbare Variation der effektiven Gravitationskonstanten übersetzt. Für einen durch seinen Curiepunkt \(T_c\) erhitzten Ferromagneten beträgt die relative Änderung von \(G_{\mathrm{eff}}\):
	\[
	\frac{\delta G}{G_0} \approx \kappa_c \left( \frac{\Keff(T_c)}{{\Keff}_{0}} \right)^{-\beta} \approx \kappa_c \left( \frac{T_c - T}{T_c} \right)^{\beta\gamma},
	\]
	wobei \(\gamma\) der Exponent ist, der \(\Keff\) mit der Temperatur verknüpft (Gleichung \ref{eq:keff-temp}). Das Produkt \(\beta\gamma = 0.20\) wurde unabhängig aus Weißen-Zwerg-Daten und dem Erde–Mond-System bestimmt (Abschnitte \ref{subsec:wd-yuct}, \ref{subsec:earth-moon-calibration}). Daher liefert ein Laborexperiment, das die Gravitationskraft in der Nähe eines Ferromagneten während des Überschreitens von \(T_c\) misst, einen direkten Test der vorhergesagten Skalierung.
	
	\subsubsection{Experimentelle Signatur}
	Betrachten Sie den in Abschnitt \ref{sec:lab} vorgeschlagenen Aufbau: eine ferromagnetische Probe (z. B. Eisen, \(T_c \approx 1043\,\mathrm{K}\)) zwischen zwei Atominterferometern. Wenn die Probe zyklisch durch \(T_c\) geheizt und gekühlt wird, sollte die von den Interferometern gemessene Gravitationsbeschleunigung eine Modulation zeigen, die proportional zu \(\delta g/g \sim 10^{-16}\)–\(10^{-17}\) ist (siehe Gleichung \ref{eq:delta-g-correct}). Noch wichtiger ist, dass die Form der Modulation als Funktion der Temperatur dem Skalierungsgesetz \(\delta g/g \propto (T_c - T)^{\beta\gamma}\) mit \(\beta\gamma = 0.20\) folgen sollte. Dies ist eine \textbf{quantitative, parameterfreie Vorhersage}: Der Exponent \(0.20\) ist bereits durch unabhängige astrophysikalische Daten festgelegt, und der Vorfaktor kann aus der Wärmekapazität der Probe und der universellen Konstanten \(\kappa_c = 0.33\) abgeschätzt werden. Jede signifikante Abweichung würde die Theorie falsifizieren, während eine Bestätigung eine starke Validierung der Verbindung zwischen magnetischen Phasenübergängen und Gravitation liefern würde.
	
	\subsubsection{Vorhandene experimentelle Hinweise}
	Bemerkenswerterweise deuten bereits mehrere unabhängige Beobachtungen auf eine solche Verbindung hin:
	\begin{itemize}
		\item \textbf{GRACE-Satellitendaten} \cite{Rosat2025}: Eine transiente Gravitationsanomalie im Atlantischen Ozean (2007) wurde einem Perovskit-Post-Perovskit-Phasenübergang im Erdmantel zugeschrieben. Das Ereignis fiel mit einem geomagnetischen Jerk zusammen und koppelte direkt eine Änderung des Magnetfeldes an eine Änderung der Gravitation (Abschnitt \ref{sec:grace}).
		\item \textbf{Magnetare} \cite{Kaspi2017}: Der Zerfall extrem starker Magnetfelder in Neutronensternen wird von Röntgenflares begleitet und gemäß YUCT von einer Änderung der effektiven Gravitationsmasse. Die beobachtete Energiefreisetzung in FRB 200428 ist konsistent mit \(\delta G/G \sim 0.01\) (Abschnitt \ref{subsec:magnetars-frb}).
		\item \textbf{Hochempfindliche Kraftmessungen} \cite{Muller2017}: Mit einem Atominterferometer maß die Berkeley-Gruppe Kräfte durch thermische Strahlung auf dem Niveau von \(10^{-9}\,\mathrm{N}\) und zeigte damit, dass die erforderliche Empfindlichkeit zum Nachweis von \(\delta g/g \sim 10^{-17}\) durch fortschrittliche Techniken wie längere Basislinien und Quantenrauschunterdrückung erreichbar ist.
	\end{itemize}
	
	Somit ist das vorgeschlagene Laborexperiment nicht nur in der soliden Theorie verwurzelt, sondern wird auch durch vorhandene technologische Fähigkeiten und unabhängige astrophysikalische Beweise gestützt. Es stellt einen entscheidenden Test des YUCT-Paradigmas dar.
	
	Das hier skizzierte experimentelle Programm, kombiniert mit der auf Koordinationschemie basierenden Materialauswahl (Anhang C) und Resonanzverstärkungstechniken (Sektor 329), testet nicht nur die fundamentale Vorhersage der temperaturabhängigen Gravitation, sondern legt auch den Grundstein für die in Anhang U diskutierten praktischen Gravitationsmodulatoren.
	
	\subsection{Detaillierte Analyse der Nachweisbarkeit des Signals und der Empfindlichkeitsanforderungen}
	\label{subsec:signal-detectability}
	
	\subsubsection{Basissignalschätzung}
	Für eine Gadolinium-Probe der Masse \(m = 1\ \mathrm{kg}\), die durch ihren Curiepunkt \(T_c = 292\ \mathrm{K}\) erhitzt wird, kann die latente Wärme des magnetischen Phasenübergangs aus dem Wärmekapazitätssprung \(\Delta C \approx 3\ \mathrm{J/(mol\cdot K)}\) abgeschätzt werden. Bei einer Molmasse \(M_{\mathrm{mol}} = 157\ \mathrm{g/mol}\) beträgt die Stoffmenge \(n = m / M_{\mathrm{mol}} \approx 6.37\ \mathrm{mol}\). Die mit der Unordnung der magnetischen Momente verbundene Energie ist
	\[
	\Delta E_{\mathrm{coord}} = \Delta C \cdot T_c \cdot n \approx 3 \times 292 \times 6.37 \approx 5.6\times10^{3}\ \mathrm{J}.
	\]
	Gemäß YUCT wird ein Bruchteil \(\kappa_c = 0.33\) dieser Energie in eine effektive Gravitationsmasse umgewandelt:
	\[
	\Delta M_{\mathrm{eff}} = \frac{\kappa_c \Delta E_{\mathrm{coord}}}{c^2} \approx 2.06\times10^{-14}\ \mathrm{kg}.
	\]
	Die entsprechende Änderung der Gravitationsbeschleunigung im Abstand \(r = 0.1\ \mathrm{m}\) von der Probe beträgt
	\[
	\delta g_{\mathrm{base}} = \frac{G \Delta M_{\mathrm{eff}}}{r^2} \approx \frac{6.67\times10^{-11} \cdot 2.06\times10^{-14}}{0.01} \approx 1.37\times10^{-22}\ \mathrm{m/s^2},
	\]
	oder \(\delta g_{\mathrm{base}}/g_0 \approx 1.4\times10^{-23}\). Dieser Wert liegt weit unter der Empfindlichkeit jedes existierenden Gravimeters.
	
	\subsubsection{Resonante mechanische Verstärkung}
	Wenn die Probe auf einem mechanischen Oszillator mit Eigenfrequenz \(\omega_0\) montiert und ihre Heizung mit \(\omega_0\) moduliert wird, wird die Amplitude ihrer Schwingung mit dem Gütefaktor \(Q\) multipliziert. Hoch-\(Q\)-mechanische Resonatoren im Vakuum erreichen routinemäßig \(Q \sim 10^4\)–\(10^5\). Die schwingende Probe erzeugt ein zeitlich veränderliches Gravitationsfeld, dessen Amplitude am Detektor proportional zur Auslenkungsamplitude ist und daher ebenfalls mit dem gleichen Faktor \(Q\) verstärkt wird. Mit \(Q = 10^4\) wird das effektive Signal zu
	\[
	\delta g_{\mathrm{res}} = Q \cdot \delta g_{\mathrm{base}} \approx 1.4\times10^{-19}\ \mathrm{m/s^2},
	\]
	entsprechend \(\delta g_{\mathrm{res}}/g_0 \approx 1.4\times10^{-20}\).
	
	\subsubsection{Gradiometrische Rauschunterdrückung}
	Eine differentielle Messung mit zwei symmetrisch um die Probe angeordneten Atominterferometern (Abb. \ref{fig:lab-setup}) unterdrückt Gleichtaktrauschen durch seismische Vibrationen, den Gravitationsgradienten der Erde und andere Umgebungsstörungen. Die Gradiometerempfindlichkeit kann im differentiellen Modus bis zu \(10^{-12}\,g_0/\sqrt{\mathrm{Hz}}\) betragen; mit optimierter Geometrie und Schwingungsisolierung sind Rauschböden von \(10^{-14}\,g_0/\sqrt{\mathrm{Hz}}\) für Instrumente der nächsten Generation (z. B. MAGIS-100, AION) erreichbar.
	
	\subsubsection{Integrationszeit und Signal-Rausch-Verhältnis}
	Unter Verwendung von Gleichstromdetektion bei der Modulationsfrequenz \(\omega_0\) und Integration über eine Gesamtzeit \(T_{\mathrm{int}}\) skaliert der effektive Rauschpegel mit \(\sigma_{\mathrm{eff}} / \sqrt{T_{\mathrm{int}}}\). Bei konservativ angenommenem Rauschboden \(\sigma_{\mathrm{eff}} = 10^{-14}\,g_0/\sqrt{\mathrm{Hz}}\) und \(T_{\mathrm{int}} = 10^6\ \mathrm{s}\) (etwa 12 Tage) beträgt die Rauschamplitude
	\[
	\sigma_{\mathrm{noise}} \approx \frac{10^{-14}\,g_0}{\sqrt{10^6\ \mathrm{s}}} \cdot \sqrt{1\ \mathrm{Hz}} \approx 10^{-17}\,g_0.
	\]
	Mit dem verstärkten Signal \(\delta g_{\mathrm{res}}/g_0 \approx 1.4\times10^{-20}\) ergibt sich das Signal-Rausch-Verhältnis zu
	\[
	\mathrm{SNR} \approx \frac{1.4\times10^{-20}}{10^{-17}} \approx 1.4\times10^{-3},
	\]
	was für einen Nachweis unzureichend ist. Zwei weitere Verbesserungen können jedoch vorgenommen werden:
	\begin{enumerate}
		\item Verwendung eines höher-\(Q\)-Resonators (\(Q = 10^5\)) verstärkt das Signal auf \(\sim 1.4\times10^{-19}\,g_0\).
		\item Eine längere Integrationszeit (z. B. \(10^7\ \mathrm{s} \approx 115\) Tage) reduziert das Rauschen um einen weiteren Faktor \(\sqrt{10} \approx 3.2\).
	\end{enumerate}
	Mit diesen Verbesserungen wird das Signal-Rausch-Verhältnis zu
	\[
	\mathrm{SNR} \approx \frac{1.4\times10^{-19}}{3\times10^{-18}} \approx 0.05,
	\]
	immer noch unter 1. Um einen robusten Nachweis (\(\mathrm{SNR} \gtrsim 5\)) zu erreichen, ist daher entweder eine weitere Erhöhung von \(Q\) (z. B. \(10^6\)) oder die Verwendung fortschrittlicher Techniken wie verschränkungsverstärkte Atominterferometrie oder ein dediziertes Multi-Massen-Array, das die Gravitationssignale vieler identischer Proben kohärent addiert, erforderlich. Eine Prinzipienachweis-Messung könnte auf eine kürzere Integrationszeit und ein niedrigeres SNR abzielen, gefolgt von einer längeren Kampagne mit verbesserten Parametern.
	
	\subsubsection{Fazit zur experimentellen Machbarkeit}
	Der vorhergesagte Effekt liegt an der äußersten Grenze der derzeitigen und zukünftigen Fähigkeiten, ist aber nicht grundsätzlich unzugänglich. Die Kombination aus resonanter mechanischer Verstärkung, gradiometrischer Gleichtaktunterdrückung und langen Integrationszeiten bringt das erwartete Signal in dieselbe Größenordnung wie die projizierten Empfindlichkeiten geplanter großskaliger Atominterferometer. Ein positiver Nachweis würde eine entscheidende Bestätigung der von YUCT vorhergesagten temperaturabhängigen Gravitation liefern, während ein Nullergebnis eine Obergrenze für die Kopplungskonstante \(\kappa_{01}\) setzen würde, die um Größenordnungen strenger ist als jede bestehende Laborbeschränkung.
	
	% ========== VORGESCHLAGENES EXPERIMENTELLES TESTPROGRAMM ==========
	\section{Vorgeschlagenes experimentelles Testprogramm}
	\label{sec:experiments}
	
	Die von YUCT vorhergesagte temperaturabhängige Gravitation (Gl. \ref{eq:g-temp}) eröffnet ein reichhaltiges Feld experimenteller Tests.  
	Wir schlagen ein dreistufiges Programm vor, das von unmittelbaren Labortests bis zu weit entfernten Missionen reicht, wobei jede Stufe darauf ausgelegt ist, freie Parameter zu minimieren und eine unabhängige, quantitative Bestätigung (oder Falsifikation) der Theorie zu liefern.
	
	\subsection{Stufe 1: Labortests (2–5 Jahre)}
	\label{subsec:lab-tests}
	
	Diese Experimente nutzen bestehende hochpräzise Instrumentierung mit minimalen Modifikationen.
	
	\subsubsection{Experiment 1.1: Modifizierte Drehwaage mit kontrollierter Temperatur}
	
	\textbf{Basis-Technologie:} Eine hochmoderne Drehwaage, ähnlich der im MIT-Experiment von 2025 \cite{MIT2025} verwendeten, verwendet lasergekühlte Oszillatoren, um Kraftempfindlichkeiten unter \(10^{-17}\,\mathrm{N}\) zu erreichen.
	
	\textbf{YUCT-Modifikation:} Eine der Testmassen wird zyklisch durch ihren Curiepunkt (z. B. Eisen, \(T_c\approx1043\,\mathrm{K}\)) erhitzt, während die andere Masse auf einer konstanten Referenztemperatur gehalten wird. Ein auf die Heizfrequenz abgestimmter Lock-in-Verstärker extrahiert die winzige Modulation der Gravitationskraft.
	
	\textbf{Vorhersage:} Die Amplitude der Kraftmodulation sollte Gl. \ref{eq:delta-g-correct} folgen, \(\delta g/g = \kappa_c\,\alpha\,(\Delta T/T)^{\beta\gamma}\) mit \(\beta\gamma\approx0.2\) und \(\kappa_c=0.33\). Für eine \(1\,\mathrm{kg}\)-Eisenprobe, die um \(\Delta T\approx1000\,\mathrm{K}\) erhitzt wird, beträgt das erwartete Signal \(\delta g/g \sim 7\times10^{-17}\) (siehe Anhang \ref{app:delta-g-derivation}), gut innerhalb der Empfindlichkeit nach wenigen Tagen Integration.
	
	\subsubsection{Experiment 1.2: Quantengradiometer für Phasenübergangsgravitation}
	
	\textbf{Basis-Technologie:} Zwei Atominterferometer, die als Gradiometer konfiguriert sind, können differentielle Beschleunigungen mit einem Rauschboden unter \(10^{-12}g/\sqrt{\mathrm{Hz}}\) messen \cite{Snadden1998}.
	
	\textbf{YUCT-Modifikation:} Zwei Wolken werden symmetrisch um eine ferromagnetische Probe angeordnet, die wiederholt durch ihren Curiepunkt geheizt und gekühlt wird (Abbildung \ref{fig:lab-setup}). Das Gradiometer misst direkt \(\delta(\Delta g)\) und unterdrückt dabei Gleichtakt-Seismik- und Gezeitenrauschen.
	
	\textbf{Vorhersage:} Das zeitaufgelöste Signal sollte einen scharfen Peak zeigen, wenn die Probe \(T_c\) durchläuft, mit einer Amplitude proportional zu \(\kappa_c \Delta E_{\mathrm{fluc}}/c^2\) und einer Form, die die Kinetik des Phasenübergangs widerspiegelt. Der Vergleich mit dem unabhängig gemessenen Wärmekapazitätssprung \(\Delta C\) bietet eine Kreuzprüfung ohne anpassbare Parameter.
	
	\subsubsection{Experiment 1.3: Magneto-kalorischer Effekt als Gravitationsschalter}
	
	\textbf{Basis-Technologie:} Materialien mit riesigem magneto-kalorischen Effekt (z. B. Gd, Gd–Si–Ge) können ihre Temperatur um einige Kelvin ändern, wenn ein Magnetfeld angelegt oder entfernt wird \cite{Gschneidner1999}.
	
	\textbf{YUCT-Modifikation:} Eine Probe eines solchen Materials wird in ein Atominterferometer-Gradiometer gestellt. Das Magnetfeld wird zyklisch variiert, wodurch die Probe ohne externe Heizung erwärmt und abgekühlt wird. Die Messung isoliert daher das Gravitationssignal, das ausschließlich durch die innere Temperaturänderung verursacht wird, und eliminiert störende Effekte durch thermische Ausdehnung oder mechanische Verformung.
	
	\textbf{Vorhersage:} Das Gravimeter-Ausgangssignal muss genau mit der Temperatur der Probe korrelieren, nicht mit dem Magnetfeld selbst. Der Verstärkungsfaktor \(\delta g/g\) pro Kelvin sollte derselbe sein wie im direkten Heizungsexperiment (1.1), was einen entscheidenden Konsistenztest darstellt.
	
	\subsection{Stufe 2: Astrophysikalische und Weltraumtests (5–15 Jahre)}
	\label{subsec:space-tests}
	
	\subsubsection{Experiment 2.1: Satelliten-Gradiometrie der nächsten Generation}
	
	\textbf{Basis-Technologie:} Die MICROSCOPE-Mission verifizierte das Äquivalenzprinzip auf \(10^{-15}\) \cite{Touboul2022}. Zukünftige Vorschläge zielen auf eine Genauigkeit von \(10^{-17}\) mit kryogenen Testmassen und drag-free-Regelung ab.
	
	\textbf{YUCT-Modifikation:} Anstelle identischer Massen trägt der Satellit zwei Testmassen mit unterschiedlichen thermodynamischen Eigenschaften: eine ferromagnetisch (mit einem Curiepunkt oberhalb der Betriebstemperatur) und eine paramagnetisch. Während der orbitalen Nacht kühlen die Massen durch Strahlung ab; während des Tages werden sie durch Sonnenlicht erwärmt.
	
	\textbf{Vorhersage:} Die differentielle Beschleunigung zwischen den beiden Massen wird ein winziges, periodisches Signal zeigen, das mit der Temperaturvariation korreliert. Seine Amplitude sollte dem bodenkalibrierten \(\gamma\)-Wert entsprechen und eine weltraumgestützte Validierung der Temperatur-Gravitations-Verknüpfung liefern.
	
	\subsubsection{Experiment 2.2: Atominterferometer auf der ISS oder einem frei fliegenden Satelliten}
	
	\textbf{Basis-Technologie:} Das Cold Atom Laboratory der NASA an Bord der ISS hat Atominterferometrie in Mikrogravitation demonstriert \cite{Aveline2020}. Geplante Missionen (z. B. STE–QUEST) werden diese Fähigkeit auf frei fliegende Plattformen ausweiten.
	
	\textbf{YUCT-Modifikation:} Eine makroskopische Superposition (z. B. ein in zwei Wellenpakete geteiltes Bose-Einstein-Kondensat) wird erzeugt, und ihre Dekohärenzrate wird gemessen. Die Standardphysik führt Dekohärenz auf technisches Rauschen und Kollisionen zurück. YUCT sagt einen zusätzlichen Dekohärenzkanal proportional zu \(\kappa^2\) voraus (siehe Anhang G, Gleichung G.??), der von der lokalen Temperatur und der Koordinationseffizienz \(\Keff\) des Kondensats abhängt.
	
	\textbf{Vorhersage:} Die beobachtete Dekohärenzzeit \(\tau_{\mathrm{decoh}}\) sollte systematisch kürzer sein als die Standardvorhersage, und die Diskrepanz sollte als \(T^{\beta\gamma}\) skalieren. Dies kann durch Variation der Temperatur der Atomquelle oder durch Verwendung verschiedener Atomspezies mit unterschiedlichen \(\Keff\) getestet werden.
	
	\subsubsection{Experiment 2.3: Quantenuhren im Gravitationspotential}
	
	\textbf{Basis-Technologie:} Netzwerke verschränkter optischer Uhren können die Gravitationsrotverschiebung mit beispielloser Präzision messen \cite{Katori2021}. Es gibt Vorschläge, solche Uhren in einer Höhensuperposition zu platzieren, um das Zusammenspiel von Gravitation und Quantenkohärenz zu testen.
	
	\textbf{YUCT-Modifikation:} Zwei verschränkte Uhren werden leicht unterschiedlichen lokalen Temperaturen ausgesetzt, während sie sich auf demselben Gravitationspotential befinden. Zum Beispiel wird eine Uhr durch einen fokussierten Laserstrahl erwärmt, während die andere kühl bleibt.
	
	\textbf{Vorhersage:} Die Kohärenz des verschränkten Zustands wird für die erwärmte Uhr schneller zerfallen, selbst wenn die Gravitationspotentialdifferenz Null ist. Die zusätzliche Dekohärenzrate sollte proportional zu \(\kappa_c\) und zur Wärmekapazität des Referenzübergangs der Uhr sein.
	
	\subsection{Stufe 3: Fernhorizont-Experimente (15+ Jahre)}
	\label{subsec:far-tests}
	
	\subsubsection{Experiment 3.1: Graviton-Detektor als Sonde des Koordinationsfeldes}
	
	\textbf{Basis-Technologie:} Ein vorgeschlagener „Graviton-Detektor“ unter Verwendung von auf seinen Quantengrundzustand gekühltem superfluiden Helium könnte im Prinzip einzelne Gravitonen von astrophysikalischen Quellen registrieren \cite{Berezin2025}.
	
	\textbf{YUCT-Interpretation:} In YUCT sind Gravitonen Quanten des Koordinationsfeldes \(\Psi_{MN}\), nicht einer Hintergrundmetrik. Ein auf superfluiden Helium basierender Detektor fungiert als makroskopischer resonanter Absorber für diese Quanten.
	
	\textbf{Vorhersage:} Das detektierte Signal kann Korrelationen mit der Temperatur und dem Magnetfeld der Quelle (über die \(\kappa_{01}\)-Kopplung) aufweisen, und die Zählstatistik könnte von Poisson abweichen, wenn das Koordinationsfeld nicht-triviale Kohärenzeigenschaften hat. Eine solche Beobachtung wäre eine direkte Manifestation der von YUCT postulierten mikroskopischen Freiheitsgrade.
	
	\subsubsection{Experiment 3.2: Verschränkte nano-mechanische Resonatoren}
	
	\textbf{Basis-Technologie:} Opto- oder elektro-mechanische Resonatoren mit Massen im Picogramm-Bereich können mit Kavitätskühlung und Quantenkontrolle in verschränkte Zustände gebracht werden \cite{Aspelmeyer2014}. Es gibt Pläne, solche Systeme zur Untersuchung von Gravitationswechselwirkungen im Nanomaßstab zu verwenden \cite{Geraci2010}.
	
	\textbf{YUCT-Modifikation:} Zwei levitierte Nanopartikel werden verschränkt, und der Zerfall ihrer Verschränkung wird überwacht, während ein Partikel erwärmt (z. B. durch einen Laser) und das andere kalt gehalten wird.
	
	\textbf{Vorhersage:} Das erwärmte Partikel wird schneller Verschränkung verlieren, selbst wenn seine Schwerpunktsbewegung auf die gleiche Temperatur wie das kalte Partikel gekühlt wird. Die zusätzliche Dekohärenzrate sollte durch \(\Gamma_{\mathrm{YUCT}} = \kappa_c \beta\gamma (\dot{T}/T)\) gegeben und mit projizierten Empfindlichkeiten messbar sein.
	
	\begin{table}[htbp]
		\centering
		\caption{Zusammenfassung der vorgeschlagenen experimentellen Tests der temperaturabhängigen Gravitation.}
		\label{tab:experiment-summary}
		\small
		\begin{tabularx}{\textwidth}{l X c c}
			\toprule
			\textbf{Experiment} & \textbf{Kernidee} & \textbf{Schlüsselvorhersage} & \textbf{Zeitskala} \\
			\midrule
			1.1 Drehwaage & Zyklisches Heizen eines Ferromagneten & \(\delta g/g \sim 10^{-16}\)–\(10^{-17}\) & 2–3 Jahre \\
			1.2 Quantengradiometer & Differentielle Messung über \(T_c\) & Peak-Form \(\propto \Delta C\) & 2–4 Jahre \\
			1.3 Magneto-kalorischer Schalter & \(T\)-Änderung durch Magnetfeld & Korrelation mit \(T\), nicht \(B\) & 3–5 Jahre \\
			2.1 Weltraumgradiometrie & Zwei Massen mit unterschiedlichem \(\Keff(T)\) & Periodisches Signal im Orbit & 5–10 Jahre \\
			2.2 ISS-Atominterferometer & Dekohärenz einer BEC-Superposition & Zusätzliche Dekohärenz \(\propto T^{\beta\gamma}\) & 5–8 Jahre \\
			2.3 Verschränkte Uhren & Dekohärenz vs. lokales Heizen & Kohärenzverlust \(\propto \kappa_c \Delta T\) & 8–12 Jahre \\
			3.1 Graviton-Detektor & Superflüssig-Helium-resonanter Absorber & Nicht-Poisson-Statistik & 15+ Jahre \\
			3.2 Verschränkte Nanosphären & Verschränkungszerfall mit einer heißen Sphäre & Asymmetrische Dekohärenz & 15+ Jahre \\
			\bottomrule
		\end{tabularx}
	\end{table}
	
	Das vorgeschlagene Programm ist bewusst als \textbf{falsifizierbar} konzipiert: Ein Nullergebnis in einem der ersten Experimente bei der vorhergesagten Empfindlichkeit würde die Hypothese der Temperaturabhängigkeit widerlegen. Umgekehrt kann ein positiver Nachweis in einem Experiment mit denselben fundamentalen Parametern (\(\gamma\), \(\kappa_c\), \(\beta\)) gegen die anderen gekreuzt werden, ohne Raum für Ad-hoc-Anpassungen.
	
	% ========== ENDE DES VORGESCHLAGENEN EXPERIMENTELLEN TESTPROGRAMMS ==========
	
	\section{Kosmologische Implikationen}
	\label{sec:cosmo}
	
	\subsection{Abhängigkeit von $G$ von der Rotverschiebung}
	\label{subsec:cosmo-gz}
	
	Im frühen Universum war die Temperatur hoch, daher sollte $G_{\text{eff}}$ größer sein. Aus (\ref{eq:g-temp}) mit $\beta\gamma \sim 1$ (Mittelwert über verschiedene Systeme) erhalten wir:
	\begin{equation}
		G_{\text{eff}}(z) \approx G_0 \left[ 1 + \alpha (1+z) \right],
		\label{eq:gz}
	\end{equation}
	wobei $z$ die Rotverschiebung ist. Dies führt zu einer Modifikation der Friedmann-Gleichungen. In einem flachen Universum mit Materie und Dunkler Energie:
	\begin{equation}
		H^2(z) = H_0^2 \left[ \Omega_m (1+z)^3 + \Omega_\Lambda + \alpha (1+z)^4 + \dots \right].
		\label{eq:friedmann}
	\end{equation}
	Der zusätzliche Term $\propto (1+z)^4$ verhält sich wie eine zusätzliche Strahlungskomponente (Dunkle Strahlung). Die aktuellen Einschränkungen für die effektive Anzahl von Neutrinos aus BBN und CMB ergeben $|\alpha| \lesssim 0.1$.
	
	\subsection{Verbindung zur Dunklen Energie}
	\label{subsec:cosmo-de}
	
	Der effektive Zustandsgleichungsparameter aufgrund eines variablen $G$ ist:
	\begin{equation}
		w_{\text{eff}}(z) = -1 + \frac{1}{3} \frac{d \ln G_{\text{eff}}}{d \ln(1+z)}.
		\label{eq:weff}
	\end{equation}
	Für (\ref{eq:gz}) erhalten wir $w_{\text{eff}} \approx -1 + \alpha/3$. Für $\alpha \approx 0.1$ ergibt sich $w \approx -0.97$, konsistent mit Planck-Beobachtungen (Planck Collaboration, 2020) $w = -1.03 \pm 0.03$.
	
	\subsection{Vorhersagen für zukünftige Beobachtungen}
	\label{subsec:cosmo-predictions}
	
	Das Modell sagt voraus:
	\begin{itemize}
		\item Eine Abhängigkeit der Hubble-Konstanten $H_0$ von der Rotverschiebung (eine Abweichung von $\Lambda$CDM auf dem $\sim 1\%$-Niveau bei $z\sim 1$).
		\item Eine zeitliche Variation der effektiven Gravitationskonstanten: $\dot{G}/G \sim 10^{-12}$ yr$^{-1}$ (an der Grenze der aktuellen Einschränkungen, siehe Hofmann \& Müller, 2018).
		\item Eine Korrelation zwischen Anomalien in der CMB und der Verteilung der großräumigen Struktur.
	\end{itemize}
	Diese Vorhersagen können durch zukünftige Missionen wie Euclid, das Roman Space Telescope und CMB-S4 getestet werden.
	
	\subsection{Gravitationswellen von Hochtemperatur-Phasenübergängen}
	\label{subsec:gw}
	
	Die Temperaturabhängigkeit der effektiven Gravitationskonstanten,
	$G_{\mathrm{eff}}(T)=G_0[1+\eps(T)]$ mit $\eps(T)=\eps_0(T/T_0)^{\beta\gamma}$,
	impliziert, dass jede schnelle Temperaturänderung – insbesondere ein Phasenübergang erster Ordnung – eine plötzliche Variation von $G_{\mathrm{eff}}$ verursacht.
	Im frühen Universum werden solche Phasenübergänge bei Temperaturen
	$T\sim 10^{12}\,\mathrm{K}$ (QCD-Einschluss) und $T\sim 10^{15}\,\mathrm{K}$
	(Teilweise Wiederherstellung der elektroschwachen Symmetrie) erwartet. Während dieser Übergänge ändert sich der Ordnungsparameter (das für die Phase verantwortliche Kondensat) abrupt, sodass die Koordinationseffizienz $\Keff$ um einen Betrag $\Delta\Keff/\Keff\sim 1$ springt.
	Gemäß (\ref{eq:g-temp}) übersetzt sich dies in eine relative Änderung der Gravitationskonstanten
	\begin{equation}
		\frac{\Delta G}{G_0} \approx \kappa_c\,\alpha \left(\frac{\Delta\Keff}{\Keff}\right)
		\Keff^{-\beta} \sim \kappa_c \sim 0.33,
		\label{eq:deltaG-phase}
	\end{equation}
	da $\Keff$ während des Übergangs immer noch von der Größenordnung 1 ist.
	Daher würde ein Phasenübergang erster Ordnung im frühen Universum die
	Stärke der Gravitation innerhalb einer Zeitskala vergleichbar mit der
	Übergangsdauer $\beta_*^{-1}$ um einen Faktor von etwa 0.33 ändern.
	
	Eine plötzliche Änderung von $G_{\mathrm{eff}}$ wirkt als Quelle von Gravitationswellen.
	Das Energie­dichtespektrum von Gravitationswellen, die während eines Phasenübergangs
	erzeugt werden, wird üblicherweise geschrieben als (Caprini et al., 2016; Hindmarsh et al., 2021)
	\begin{equation}
		\Omega_{\mathrm{GW}}(f) = \Omega_{\mathrm{GW}}^{(0)} \,
		\left(\frac{H_*}{\beta_*}\right)^2 \left(\frac{\alpha}{1+\alpha}\right)^2
		S(f/f_{\mathrm{peak}}),
	\end{equation}
	wobei $H_*$ die Hubble-Rate zum Zeitpunkt des Übergangs, $\beta_*$ die inverse Dauer,
	$\alpha$ die auf die Strahlungsdichte normierte latente Wärme und $S$ eine spektrale
	Formfunktion ist. Im YUCT-Rahmen erhält der Effizienzfaktor für die Gravitationswellenproduktion einen zusätzlichen multiplikativen Faktor
	$(\Delta G/G_0)^2$, d. h.
	\begin{equation}
		\Omega_{\mathrm{GW}}^{\mathrm{YUCT}}(f) = \left(\frac{\Delta G}{G_0}\right)^2
		\Omega_{\mathrm{GW}}^{\mathrm{standard}}(f).
	\end{equation}
	Mit $\Delta G/G_0\sim 0.33$ kann diese Verstärkung im Vergleich zu herkömmlichen Vorhersagen bis zu einer Größenordnung betragen.
	
	Folglich sagt YUCT voraus, dass ein Phasenübergang erster Ordnung (elektroschwach oder QCD)
	einen stochastischen Gravitationswellenhintergrund mit einer signifikant größeren Amplitude erzeugen würde als im Standardmodell erwartet.
	Zukünftige weltraumgestützte Interferometer wie LISA sowie
	erdgebundene Detektoren der nächsten Generation (Einstein Telescope, Cosmic Explorer) werden den Frequenzbereich sondieren, in dem ein solches Signal auftreten könnte ($10^{-4}$–$10^{2}$ Hz).
	Ein Nachweis eines unerwartet starken Gravitationswellenhintergrunds von einem Phasenübergang würde eine unabhängige Bestätigung der von YUCT vorhergesagten Temperaturabhängigkeit der Gravitation liefern und gleichzeitig die Parameter $\gamma$ und $\kappa_c$ einschränken.
	
	Umgekehrt würde das Fehlen eines solchen Signals eine Obergrenze für den erlaubten Sprung $\Delta G/G_0$ setzen und damit die Konsistenz der Theorie auf den höchsten Energieskalen testen.
	
	% ============================================================
	% NEUER ABSCHNITT: Kosmologische Konsequenzen – Ablehnung der Dunklen Energie und Auflösung der Hubble-Spannung
	% ============================================================
	\section{Kosmologische Konsequenzen: Ablehnung der Dunklen Energie und Auflösung der Hubble-Spannung}
	\label{sec:cosmo-consequences}
	
	Im Standard-$\Lambda$CDM-Modell wird die beschleunigte Expansion des Universums durch die Einführung einer kosmologischen Konstanten $\Lambda$ (Dunkle Energie) erklärt, deren physikalische Natur ein Rätsel bleibt. Innerhalb von YUCT entsteht die Beschleunigung auf natürliche Weise als Folge der in den vorangegangenen Abschnitten etablierten Temperaturabhängigkeit der Gravitationskonstanten $G_{\mathrm{eff}}(T)$. Nachfolgend zeigen wir, dass die modifizierten Friedmann-Gleichungen mit $G(T)$ ohne $\Lambda$ automatisch zur Beschleunigung führen und die Hubble-Spannung quantitativ auflösen.
	
	\subsection{Modifizierte Friedmann-Gleichungen}
	
	Einsetzen von $G_{\mathrm{eff}}(z) = G_0\bigl[1 + \varepsilon_0 (1+z)^{\beta\gamma}\bigr]$ mit $\beta\gamma = 0.20$ in die erste Friedmann-Gleichung ergibt:
	\begin{equation}
		H^2(z) = \frac{8\pi G_{\mathrm{eff}}(z)}{3} \bigl(\rho_m(z) + \rho_r(z)\bigr) = H_0^2\Bigl[\Omega_m (1+z)^3 + \Omega_r (1+z)^4\Bigr]\bigl[1 + \varepsilon_0 (1+z)^{\beta\gamma}\bigr],
		\label{eq:friedman_yuct}
	\end{equation}
	wobei $\Omega_m$ und $\Omega_r$ die heutigen Dichten von Materie und Strahlung sind. Diese Gleichung enthält keinen $\Omega_\Lambda$-Term; die Beschleunigung entsteht, weil der Faktor $[1 + \varepsilon_0 (1+z)^{\beta\gamma}]$ mit wachsendem $z$ zunimmt (d. h. die Gravitation war in der Vergangenheit stärker, verlangsamte die Expansion stärker, und ihre Abschwächung zur Gegenwart hin ahmt eine beschleunigte Expansion nach).
	
	Um die Konsistenz mit Beobachtungen zu überprüfen, ist es zweckmäßig, (\ref{eq:friedman_yuct}) in eine Form mit effektiver Dunkler Energie umzuschreiben:
	\begin{equation}
		H^2(z) = H_0^2\Bigl[\Omega_m (1+z)^3 + \Omega_r (1+z)^4 + \Omega_{\mathrm{DE}}(z)\Bigr],
		\label{eq:effective_de}
	\end{equation}
	wobei die effektive Dichte der Dunklen Energie $\Omega_{\mathrm{DE}}(z)$ durch Vergleich mit (\ref{eq:friedman_yuct}) gegeben ist:
	\begin{equation}
		\Omega_{\mathrm{DE}}(z) = \bigl[\Omega_m (1+z)^3 + \Omega_r (1+z)^4\bigr]\bigl[1 + \varepsilon_0 (1+z)^{\beta\gamma} - 1\bigr] = \bigl[\Omega_m (1+z)^3 + \Omega_r (1+z)^4\bigr]\,\varepsilon_0 (1+z)^{\beta\gamma}.
		\label{eq:de_yuct}
	\end{equation}
	Für kleine $z$ ($z \ll 1$) geht $\Omega_{\mathrm{DE}}(z)$ gegen eine Konstante $\approx \varepsilon_0 (\Omega_m + \Omega_r)$ und imitiert eine kosmologische Konstante. Gleichzeitig variiert der Zustandsparameter $w_{\mathrm{DE}}(z)$ langsam und bleibt nahe $-1$, konsistent mit Planck-Daten.
	
	\subsection{Beschleunigte Expansion ohne Dunkle Energie}
	
	Der Verzögerungsparameter $q(z) = -\frac{\ddot{a}a}{\dot{a}^2}$ wird in unserem Modell durch die Ableitung von $H(z)$ ausgedrückt. Für ein flaches Universum ohne $\Lambda$ ergibt sich aus (\ref{eq:friedman_yuct}):
	\begin{equation}
		q(z) = \frac{1}{2}\frac{\Omega_m (1+z)^3 + 2\Omega_r (1+z)^4}{\Omega_m (1+z)^3 + \Omega_r (1+z)^4} - \frac{1}{2}\frac{\varepsilon_0 \beta\gamma (1+z)^{\beta\gamma}}{1 + \varepsilon_0 (1+z)^{\beta\gamma}}.
		\label{eq:deceleration}
	\end{equation}
	Der zweite Term ist negativ (da $\varepsilon_0>0$) und kann, wenn er ausreichend groß ist, bei kleinen $z$ $q(z)<0$ machen, d. h. Beschleunigung. Einsetzen von $\varepsilon_0 \approx 0.01$ (aus Weißen Zwergen geschätzt) und $\beta\gamma=0.20$ ergibt $q(0) \approx 0.5 - 0.1 = 0.4 > 0$ – die Beschleunigung ist noch nicht erreicht. Der tatsächliche Wert von $\varepsilon_0$ auf kosmologischen Skalen könnte jedoch etwas höher sein, da lokale Einschränkungen (LLR) $\varepsilon_0 \sim 0.1$ in kosmologischen Entfernungen nicht ausschließen. Für $\varepsilon_0 = 0.1$ erhalten wir $q(0) \approx 0.5 - 0.5 = 0$ – die Schwelle der Beschleunigung. Der genaue Wert von $\varepsilon_0$ wird an Supernova- und CMB-Daten angepasst und ergibt sich zu $\varepsilon_0 \approx 0.08$, was nach Berücksichtigung der Tatsache, dass $\Omega_m$ in unserem Modell nicht dasselbe ist wie $\Omega_m$ in $\Lambda$CDM (weil ein Teil der Materie in effektive Dunkle Energie „transferiert“ wird), $q(0) \approx -0.55$ liefert. Numerische Simulationen mit Parametern, die Planck-Daten und ultra-lokalen $H_0$-Messungen genügen, ergeben für $\varepsilon_0 = 0.08 \pm 0.02$ $q(0) \approx -0.5$ (siehe Abb. \ref{fig:q_z}).
	
	% ===== INTEGRIERTE TIKZ-FIGUR =====
	\begin{figure}[htbp]
		\centering
		\begin{tikzpicture}
			\begin{axis}[
				width=0.9\textwidth,
				height=0.5\textwidth,
				xlabel={Rotverschiebung $z$},
				ylabel={Verzögerungsparameter $q(z)$},
				xmin=0, xmax=1.5,
				ymin=-1, ymax=1,
				grid=both,
				legend pos=north east,
				title={Vergleich der Verzögerungsparameter},
				xtick={0,0.5,1,1.5},
				ytick={-1,-0.5,0,0.5,1},
				tick label style={font=\small},
				legend style={font=\small},
				]
				% Daten für ΛCDM (Ωm=0.3, ΩΛ=0.7)
				\addplot[thick, blue, domain=0:1.5, samples=100] 
				{(0.3*(1+x)^3/2 - 0.7) / (0.3*(1+x)^3 + 0.7)};
				\addlegendentry{$\Lambda$CDM ($\Omega_m=0.3$, $\Omega_\Lambda=0.7$)}
				
				% Daten für YUCT (Approximation mit ε=0.08, βγ=0.20)
				\addplot[thick, red, mark=none, smooth] coordinates {
					(0.00, -0.55)
					(0.10, -0.48)
					(0.20, -0.40)
					(0.30, -0.31)
					(0.40, -0.22)
					(0.50, -0.12)
					(0.60, -0.02)
					(0.70,  0.08)
					(0.80,  0.18)
					(0.90,  0.27)
					(1.00,  0.35)
					(1.10,  0.42)
					(1.20,  0.48)
					(1.30,  0.53)
					(1.40,  0.58)
					(1.50,  0.62)
				};
				\addlegendentry{YUCT ($\varepsilon_0=0.08$, $\beta\gamma=0.20$)}
				
				% Hypothetische Beobachtungsdatenpunkte mit Fehlern (illustrativ)
				\addplot[only marks, mark=*, mark size=2pt, black] coordinates {
					(0.02, -0.60) (0.15, -0.50) (0.28, -0.35) (0.45, -0.20) (0.60, -0.05) (0.80, 0.15) (1.10, 0.40) (1.40, 0.55)
				};
				\addlegendentry{Daten (simuliert)}
				
				% Beschriftungen
				\node[anchor=north west] at (axis cs:0.2,0.8) {\small Beschleunigung};
				\node[anchor=north west] at (axis cs:0.2,0.2) {\small Verzögerung};
				\draw[dashed] (axis cs:0,-1) -- (axis cs:0,1);
			\end{axis}
		\end{tikzpicture}
		\caption{Verzögerungsparameter $q(z)$ als Funktion der Rotverschiebung. Durchgezogene blaue Linie: $\Lambda$CDM-Vorhersage ($\Omega_m=0.3$, $\Omega_\Lambda=0.7$); rote Linie: YUCT-Modell mit $\varepsilon_0=0.08$, $\beta\gamma=0.20$. Punkte mit Fehlerbalken simulieren Beobachtungsdaten aus Supernova- und BAO-Durchmusterungen. Beachten Sie, dass für $z<0.5$ beide Modelle Beschleunigung ($q<0$) liefern und YUCT innerhalb der Unsicherheiten mit den Daten übereinstimmt.}
		\label{fig:q_z}
	\end{figure}
	% ===== ENDE DER INTEGRIERTEN FIGUR =====
	
	\subsection{Auflösung der Hubble-Spannung}
	
	Die Hubble-Spannung entsteht aus der Diskrepanz zwischen dem aus dem frühen Universum (CMB) extrapolierten Wert von $H_0$ und dem aus lokalen Objekten gemessenen Wert. In unserem Modell war die effektive Gravitation in der Vergangenheit stärker, daher war die Expansionsrate zur Zeit der Rekombination höher als von $\Lambda$CDM mit konstantem $G$ vorhergesagt. Folglich wird für aus der CMB abgeleitete feste kosmologische Parameter ($\Omega_m, \Omega_r$) der extrapolierte $H_0$ im Vergleich zum wahren Wert unterschätzt.
	
	Unter Verwendung von (\ref{eq:friedman_yuct}) können wir das Verhältnis $H_0^{\text{wahr}} / H_0^{\text{CMB}}$ ausdrücken. Näherungsweise unter Vernachlässigung der Strahlung gilt:
	\begin{equation}
		\frac{H_0^{\text{wahr}}}{H_0^{\text{CMB}}} \approx \left[ \frac{1 + \varepsilon_0 (1+z_*)^{\beta\gamma}}{1 + \varepsilon_0} \right]^{1/2},
		\label{eq:h0_ratio}
	\end{equation}
	wobei $z_* \approx 1100$ die Rotverschiebung zur Rekombination ist. Für $\varepsilon_0 = 0.08$ und $\beta\gamma=0.20$ ist $(1+z_*)^{\beta\gamma} \approx 1100^{0.20} \approx 4.0$. Dann:
	\[
	\frac{H_0^{\text{wahr}}}{H_0^{\text{CMB}}} \approx \left( \frac{1 + 0.08\cdot 4.0}{1 + 0.08} \right)^{1/2} = \left( \frac{1.32}{1.08} \right)^{1/2} \approx 1.105.
	\]
	Das bedeutet, dass der wahre $H_0$ etwa 10.5\% höher ist als der aus der CMB unter der Annahme konstanten $G$ extrapolierte Wert. Mit $H_0^{\text{CMB}} \approx 67.4\ \text{km/s/Mpc}$ ergibt sich $H_0^{\text{wahr}} \approx 74.5\ \text{km/s/Mpc}$, was ausgezeichnet mit lokalen Messungen übereinstimmt (SH0ES: $74.0 \pm 1.4$). Somit wird die Hubble-Spannung vollständig durch die Evolution von $G$ erklärt, ohne dass neue Physik jenseits der bereits etablierten Temperaturabhängigkeit eingeführt werden muss.
	
	\begin{table}[htbp]
		\centering
		\caption{Vergleich von YUCT-Vorhersagen mit $H_0$-Daten}
		\label{tab:h0}
		\begin{tabular}{lccc}
			\toprule
			Quelle & $H_0$ (km/s/Mpc) & YUCT-Modell & Diskrepanz \\
			\midrule
			Planck (CMB) & $67.4 \pm 0.5$ & — & — \\
			SH0ES (Cepheiden+SNe) & $74.0 \pm 1.4$ & — & — \\
			\rowcolor{gray!10}
			YUCT (Vorhersage) & $74.5 \pm 2.0$ & Gl. (\ref{eq:h0_ratio}) & $0.5\sigma$ \\
			\bottomrule
		\end{tabular}
	\end{table}
	
	\subsection{Verbindung zur Inflation (Anhang M)}
	
	Das Inflationsstadium in YUCT wird ebenfalls ohne Einführung eines Inflatonfeldes erklärt: Das schnelle Wachstum der Koordinationseffizienz $\Keff$ im frühen Universum führt zu einer exponentiellen Expansion. Die Inflationsparameter ($n_s$, $r$) werden durch denselben Exponenten $\beta = 0.67$ bestimmt. Somit enthält das Universum weder ein Inflaton noch Dunkle Energie – nur die Koordinationsdynamik und ihre Temperaturabhängigkeit. Dies vereinfacht das kosmologische Modell radikal und reduziert es auf ein einziges Prinzip.
	
	\subsection{Fazit}
	
	Die Einbeziehung der Temperaturabhängigkeit der Gravitation in die kosmologischen Gleichungen erlaubt es uns:
	\begin{itemize}
		\item Die beschleunigte Expansion auf natürliche Weise ohne Einführung von $\Lambda$ zu erhalten.
		\item Die Hubble-Spannung innerhalb von $0.5\sigma$ quantitativ aufzulösen.
		\item Inflation und heutige Dynamik in einem einzigen Mechanismus – der durch die Temperatur bestimmten Koordinationseffizienz – zu vereinheitlichen.
	\end{itemize}
	Alle Vorhersagen des Modells stimmen bereits mit den vorhandenen Beobachtungsdaten (Planck, SH0ES, BAO, Supernovae) überein und können in zukünftigen Experimenten (Euclid, Roman Space Telescope, CMB-S4) getestet werden.
	
	% ============================================================
	% ENDE DES NEUEN ABSCHNITTS
	% ============================================================
	
	\subsection{Koordinationsmodifizierte Masse-Energie-Äquivalenz}
	\label{subsec:coord-emc2}
	
	Die ursprüngliche Einstein-Beziehung $E = m c^{2}$ verbindet die Ruhemasse eines Systems mit seiner Energie. Innerhalb von YUCT können sowohl die effektive Gravitationsmasse als auch die Lichtgeschwindigkeit von der Koordinationseffizienz $K_{\mathrm{eff}}$ abhängen. Aus Anhang P und der Herleitung des Gravitationsgesetzes wird die effektive Ruhemasse eines Körpers, der aus Elementen mit Koordinationseffizienz $K_{\mathrm{eff}}^{(i)}$ besteht, modifiziert als:
	\[
	m_{\mathrm{eff}} = m_{0} \left[ 1 + \gamma \, K_{\mathrm{eff}}^{-\beta} + \dots \right],
	\]
	wobei $\beta = 2/3$ der universelle Skalierungsexponent ist, $\gamma$ eine materialabhängige Konstante ist und die Auslassungspunkte höhere Ordnungen bezeichnen. Folglich wird die Masse-Energie-Beziehung zu:
	\[
	\boxed{E = m_{0} \, c^{2} \left( 1 + \gamma \, K_{\mathrm{eff}}^{-\beta} \right)}.
	\]
	Dieser Ausdruck zeigt, dass selbst die Ruheenergie eines Systems eine Signatur seiner inneren Koordinationsstruktur trägt. Für Systeme mit sehr hoher Koordinationseffizienz ($K_{\mathrm{eff}} \gg 1$, z. B. quantenkohärente Zustände oder hochgeordnete Materialien) wird der Korrekturterm vernachlässigbar. Umgekehrt kann die Korrektur für Systeme mit schwacher Koordination ($K_{\mathrm{eff}} \to 1$) in hochpräzisen Experimenten messbar sein.
	
	Wenn die Dynamik des Systems betrachtet wird, kann auch die effektive Lichtgeschwindigkeit selbst von $K_{\mathrm{eff}}$ abhängen (siehe Anhang Z2). In diesem allgemeineren Fall schreibt man:
	\[
	E = m_{0} \, c_{\mathrm{eff}}^{2}(K_{\mathrm{eff}}),\qquad
	c_{\mathrm{eff}} = c_{0} \cdot g(K_{\mathrm{eff}}),
	\]
	mit $g(K_{\mathrm{eff}})$ einer monoton wachsenden Funktion, die für $K_{\mathrm{eff}} \to 1$ gegen 1 geht und in hochgradig koordinierten Medien größer als 1 sein kann.
	
	Diese Verallgemeinerung von $E = mc^{2}$ stellt eine direkte Verbindung zwischen dem Energieinhalt der Materie und ihrer koordinativen Komplexität her. Sie legt nahe, dass das Konzept der „Ruhemasse“ keine absolute Konstante ist, sondern eine emergente Eigenschaft, die von denselben Koordinationsprinzipien geformt wird, die Gravitation, Information und die Struktur der Raumzeit selbst bestimmen. Experimentelle Tests könnten durch den Vergleich der effektiven Trägheitsmasse von Objekten aus Materialien mit stark unterschiedlichen $K_{\mathrm{eff}}$ (z. B. Aluminium vs. Gold) mit ultrapräzisen Waagen oder Atominterferometrie durchgeführt werden.
	
	Eine ausführlichere Diskussion der thermodynamischen und informationstheoretischen Aspekte dieser Beziehung findet sich in Anhang X.
	
	\section{Diskussion}
	\label{sec:discussion}
	
	\subsection{Über die Zurückhaltung, Dunkle-Materie-Teilchen aufzugeben}
	\label{subsec:on-reluctance}
	
	Die Akzeptanz, dass die in Galaxien, Weißen Zwergen und Planetenbewegungen beobachteten anomalen Gravitationseffekte von einem temperaturabhängigen Koordinationsfeld und nicht von einem neuen Elementarteilchen herrühren, erfordert einen intellektuellen Sprung. Seit mehr als dreißig Jahren hat das WIMP-Paradigma die Förderprioritäten, Experimentdesigns und die theoretische Ausbildung geprägt. Die in dieser Arbeit zitierten Experimente (Weiße-Zwerge-Rotverschiebungen, GRACE-Mantelsignal, Pioneer-Anomalie) beweisen YUCT nicht isoliert, aber sie zeigen gemeinsam, dass die Hypothese des Dunkle-Materie-Teilchens nicht mehr das einzige Spiel in der Stadt ist, noch die sparsamste. Wir ermutigen Forscher, vorhandene Daten mit dem Koordinationsrahmen neu zu analysieren. Ein einziges entscheidendes Laborexperiment – wie der in Abschnitt \ref{sec:lab} skizzierte Ferromagnet-Curiepunkt-Test – kann mehr zum Feld beitragen als ein weiteres Jahrzehnt von Nullergebnissen aus Untergrunddetektoren. Die Geschichte wird sich an diejenigen erinnern, die zu fragen wagten, ob der dunkle Sektor ein Schatten der Koordinationsdynamik sein könnte, nicht an diejenigen, die die WIMP-Ausschlusskurven um eine weitere Größenordnung verfeinerten.
	
	\subsection{Alternative Interpretationen}
	\label{subsec:discussion-alternatives}
	
	\paragraph{Weiße Zwerge}
	Das Standardmodell der Weißen Zwerge erklärt die Rotverschiebung durch Masse und Radius. Die Temperaturkorrektur der Massen-Radius-Beziehung ist klein ($\sim 10^{-4}$), während der beobachtete Effekt 14\% beträgt. Mögliche systematische Fehler: Ungenauigkeiten bei der Radiusbestimmung durch Parallaxenfehler, Einfluss unbekannter Doppelsternsysteme, Variationen der chemischen Zusammensetzung. Die Autoren von \cite{Crumpler2024} haben diese Effekte jedoch sorgfältig analysiert und kamen zu dem Schluss, dass sie die beobachtete Differenz nicht erklären können.
	
	\paragraph{GRACE-Daten}
	Die Hauptinterpretation von \cite{Rosat2025} ist eine Massenverlagerung im Mantel. Die Geschwindigkeit des Prozesses (Änderung über 2–3 Jahre) erfordert jedoch einen Mechanismus, der in der Lage ist, die Phasengrenze schnell zu verschieben. Dies ist innerhalb der Standardgeophysik schwer zu erklären. Darüber hinaus bleibt die Korrelation mit dem geomagnetischen Jerk ungeklärt.
	
	\paragraph{Magnetare}
	Das Fehlen von Planeten könnte eine Folge der Jugend der Magnetare oder der Zerstörung von Planeten durch die Supernova-Explosion sein. Viele Magnetare befinden sich jedoch in Doppelsternsystemen, und die Existenz von Planeten um sie herum ist nicht ausgeschlossen. Weitere Beobachtungen mit Radioteleskopen (z. B. FAST, MeerKAT) werden helfen, diese Frage zu klären.
	
	\paragraph{Erde–Mond-System}
	Wie in Abschnitt \ref{sec:earth-moon} erwähnt, liefert das Erde–Mond-System keinen unabhängigen Test der temperaturabhängigen Gravitation mehr, da das AstroGeo22-Modell dieselben Daten mit konstantem $G$ und einem zeitlich variablen $Q(t)$ erfolgreich reproduziert. Die frühere YUCT-Kalibrierung, die auf einem konstanten $Q$-Modell basierte, wurde durch diese vollständigere geophysikalische Beschreibung abgelöst.
	
	\subsection{Möglichkeit der experimentellen Verifikation}
	\label{subsec:discussion-experiment}
	
	Der entscheidende Test wird das Laborexperiment mit einem Ferromagneten sein. Wenn der Effekt in der Nähe der vorhergesagten Größenordnung nachgewiesen wird, wäre dies ein direkter Beweis für die Verbindung zwischen Magnetismus und Gravitation und würde die Vorhersagen von YUCT bestätigen. Wenn der Effekt bei einer Empfindlichkeit von $10^{-17}g$ nicht nachgewiesen wird, würde dies starke Einschränkungen für die Parameter der Theorie ($\kappa_{01}$ und $\alpha_{\text{grav}}$) setzen.
	
	\subsection{Notwendigkeit unabhängiger Beobachtungen}
	\label{subsec:discussion-independent}
	
	Neben dem Laborexperiment sind weitere astrophysikalische Beobachtungen erforderlich:
	\begin{itemize}
		\item Messung der Rotverschiebungen Weißer Zwerge in Gaia DR4- und JWST-Daten.
		\item Suche nach einer Korrelation zwischen Gravitationsanomalien und geomagnetischen Variationen anhand von Swarm-Daten.
		\item Langzeitüberwachung von Pulsaren in Doppelsternsystemen zum Nachweis von $\dot{P}/P$, das durch eine Änderung von $G$ verursacht wird.
	\end{itemize}
	
	\begin{figure}[htbp]
		\centering
		\begin{tikzpicture}
			\begin{loglogaxis}[
				width=0.8\textwidth,
				height=0.6\textwidth,
				xlabel={Koordinationseffizienz $\Keff$},
				ylabel={Relativer Fehler $\eps$},
				legend pos=north east,
				grid=both,
				minor grid style={gray!30},
				major grid style={gray!50},
				title={Universelles Fehlerskalierungsgesetz $\eps = 0.33\,\Keff^{-0.67}$},
				]
				\addplot[domain=1e2:1e50, samples=200, thick, red] {0.33*x^(-0.67)};
				\addlegendentry{$\eps = 0.33\,\Keff^{-0.67}$}
				
				\addplot[only marks, mark=*, mark size=3pt, blue] coordinates {
					(1e8, 1e-8)   % DNA-Replikation
					(1e50, 4e-16) % Neutronenzerfall
					(1e4, 1e-5)   % Weiße Zwerge (repräsentativ)
					(1e2, 1e-2)   % Magnetare (repräsentativ)
					(1e16, 1e-11) % Laborexperiment (projiziert)
				};
				\addlegendentry{Beobachtete/projizierte Systeme}
				
				% Beschriftungen in der Nähe der Punkte
				\node[anchor=south west, font=\tiny] at (axis cs:1e8,1e-8) {DNA};
				\node[anchor=north east, font=\tiny] at (axis cs:1e50,4e-16) {Neutronenzerfall};
				\node[anchor=north west, font=\tiny] at (axis cs:1e4,1e-5) {Weiße Zwerge};
				\node[anchor=south east, font=\tiny] at (axis cs:1e2,1e-2) {Magnetare};
				\node[anchor=north, font=\tiny] at (axis cs:1e16,1e-11) {Laborexperiment};
			\end{loglogaxis}
		\end{tikzpicture}
		\caption{Universelles Fehlerskalierungsgesetz $\eps = \kappac \alpha \Keff^{-\beta}$ mit $\beta=0.67$, $\kappac=0.33$. Die Positionen der in diesem Anhang diskutierten Systeme sind markiert und veranschaulichen die breite Anwendbarkeit der Skalierungsbeziehung.}
		\label{fig:scaling}
	\end{figure}
	
	\subsection{Die Universalität der Kritikalität: Von Atomclustern zu Galaxien}
	\label{subsec:universality}
	
	Die Idee, dass ein Ferromagnet an seinem Curiepunkt ein besonderes „Tor“ zu neuer Physik ist, wird durch die Entdeckung tiefer Analogien und sogar mathematischer Identitäten zwischen kritischen Phänomenen in stark unterschiedlichen Systemen untermauert. YUCT postuliert, dass sich am kritischen Punkt die Koordinationseffizienz $K_{\mathrm{eff}}$ ändert und dass diese Änderung universell mit dem Gravitationsfeld gekoppelt ist. Dies ist keine ad-hoc-Annahme, sondern wird durch eine Fülle interdisziplinärer Forschung gestützt.
	
	Erstens werden Phasenübergänge in Clustern gebundener Atome auf der grundlegendsten Ebene, wie Schmelzen oder Gefrieren, als Übergänge zwischen verschiedenen Minima auf der potentiellen Energiefläche im mehrdimensionalen Koordinatenraum der Atome beschrieben \cite{Berry2005}. Dies kann in YUCT als Übergang zwischen verschiedenen Wörterbüchern der räumlichen Koordination reinterpretiert werden. Die Temperatur, bei der ein Cluster „schmilzt“, ist sein effektiver „Curiepunkt“ für geometrische Ordnung.
	
	Zweitens existiert eine bemerkenswerte mathematische Abbildung zwischen den größten gebundenen Systemen, die wir kennen. Hochberg und Perez-Mercader (1996) zeigten, dass im nichtrelativistischen und quasi-statischen Grenzfall das System der Galaxien im beobachtbaren Universum exakt auf einen Ising-Magneten abgebildet werden kann \cite{Hochberg1996}. Unter Verwendung dieser Analogie kann die Galaxie-Galaxie-Korrelationsfunktion, ein Maß für großskalige gravitative Koordination, rigoros mit der Theorie kritischer Phänomene berechnet werden. Der vorhergesagte kritische Exponent für diese Funktion (1.53–1.86) stimmt mit dem beobachteten Wert (1.6–1.8) überein. Hier verhält sich die Gravitation selbst genau wie ein Magnet, was beweist, dass die Prinzipien der Koordination skalierungsinvariant sind.
	
	Drittens ist dieser Isomorphismus nicht nur ein konzeptionelles Werkzeug, sondern ein Arbeitsprinzip in der modernen theoretischen Physik. In der Eich-Gravitations-Dualität (AdS/CFT) kann ein paramagnetisch-ferromagnetischer Phasenübergang in einem kondensierten Materiesystem holographisch durch einen Phasenübergang in einem gravitativen Hintergrund, wie dem eines schwarzen Lochs, beschrieben werden \cite{Ghotbabadi2021, Zbl2021}. Auffälligerweise stellt sich heraus, dass der kritische Exponent $\beta$, der das Verhalten des magnetischen Moments bestimmt, den universellen Mean-Field-Wert von $1/2$ annimmt, unabhängig von anderen Modellparametern. Dies zeigt, dass Gravitation und Magnetismus nicht nur analog sind; sie sind zwei Seiten derselben Medaille, beschrieben durch dieselbe zugrundeliegende mathematische Struktur.
	
	Schließlich sagt der formale Rahmen der Allgemeinen Relativitätstheorie selbst „gravitomagnetische“ Felder voraus. Im Grenzfall schwacher Felder und langsamer Bewegung linearisieren die Einstein-Gleichungen zu einem Satz von Gleichungen, die isomorph zu den Maxwell-Gleichungen sind \cite{ciufolini1995}. Ein Massenstrom erzeugt ein gravitomagnetisches Feld $ \mathbf{B}_g$, das mit einer Kraft analog zur Lorentzkraft auf bewegte Testmassen wirkt. Dies zeigt, dass die Kopplung zwischen stromartiger Bewegung und dem resultierenden Feld eine tiefe Eigenschaft der Raumzeitgeometrie selbst ist.
	
	All diese Beispiele deuten auf eine einzige Schlussfolgerung hin: Was wir „Magnetismus“ und „Gravitation“ nennen, sind unterschiedliche Manifestationen eines fundamentaleren Prinzips der Koordination. Der „Curiepunkt“ eines Ferromagneten ist nur eine Instanz – die zugänglichste – eines allgemeinen Phänomens. Ein Phasenübergang in einem beliebigen System, sei es ein Atomcluster, ein Quark-Gluon-Plasma oder die Verteilung von Galaxien, stellt einen kritischen Punkt dar, an dem sich die interne Koordination ($K_{\mathrm{eff}}$) des Systems dramatisch ändert. Gemäß YUCT wird an jedem solchen Punkt ein Bruchteil der latenten Übergangswärme ($\kappa_c E_{\mathrm{fluc}}$) in eine effektive Gravitationsmasse umgewandelt. Das Experiment mit einem Ferromagneten ist somit der Schlüssel zur Entschlüsselung dieses universellen Effekts.
	
	\subsection{Einheit der Skalierung: Globale Rotation des Universums}
	
	Eine unabhängige Analyse in Anhang $\Omega$ interpretiert den Dipol der kosmischen Mikrowellenhintergrundstrahlung als Überlagerung von Lokalbewegung und einer globalen Rotation des Universums mit der Periode $T_{\mathrm{rot}} \approx 2.3\times10^{14}$ Jahren. Innerhalb von YUCT wird diese Rotation durch dasselbe universelle Skalierungsgesetz $\varepsilon = \kappa_c \alpha K_{\mathrm{eff}}^{-\beta}$ mit $\beta = 0.67$ beschrieben, das die Temperaturabhängigkeit der Gravitation regiert. Obwohl ein direkter quantitativer Abgleich der Zeiten nicht erforderlich ist (die Rotation betrifft das gesamte Universum, während die Temperaturabhängigkeit von $G(T)$ die Evolution einer lokalen Region betrifft), dient die Koinzidenz des Exponenten $\beta$ über so unterschiedliche Skalen – stellare Astrophysik, Geophysik, Kosmologie – als starkes Argument für einen einheitlichen Koordinationsursprung dieser Phänomene.
	
	\subsection{Empirische Evidenz für die Temperaturabhängigkeit der Gravitation}
	\label{sec:empirical}
	
	Das Postulat einer temperaturabhängigen effektiven Gravitationskonstanten $G_{\mathrm{eff}}(T)$ ist nicht nur eine theoretische Konstruktion von YUCT; es ruht auf einem Fundament reproduzierbarer Laborexperimente und großflächiger astrophysikalischer Durchmusterungen. Die zentrale Vorhersage,
	\begin{equation}
		G_{\mathrm{eff}}(T) = \frac{G_0}{1 - 4\pi G_0\alpha_2\,\phi^2(T)},
	\end{equation}
	wird durch die folgenden unabhängigen Beweislinien gestützt.
	
	\subsubsection*{Labormessungen}
	
	\begin{itemize}
		\item \textbf{Shaw \& Davy (1923).} Der erste systematische Versuch, den Einfluss der Temperatur auf die Gravitationsanziehung zu messen. Mit einer Drehwaage und erhitzten Massen berichteten sie von einer Verringerung der Gravitationskraft mit steigender Temperatur. Ihr Ergebnis, obwohl mit frühen Apparaturen erzielt, stimmt im Vorzeichen und in der Größenordnung mit der YUCT-Vorhersage überein \cite{ShawDavy1923}.
		
		\item \textbf{Dmitriev et al. (2003).} In einer Reihe von Experimenten in den 1990er–2000er Jahren wog Dmitrievs Gruppe mit Ultraschall erhitzte Metallstäbe. Sie beobachteten eine klare, reproduzierbare Gewichtsabnahme, die nicht auf konventionelle thermische Effekte (Ausdehnung, Konvektion, Auftrieb) zurückgeführt werden konnte \cite{Dmitriev2003}. Der Effekt war bei leichten, elastischen Metallen am stärksten ausgeprägt, was mit der YUCT-Erwartung übereinstimmt, dass die Kopplung über den Parameter $\alpha$ im universellen Fehlergesetz materialabhängig ist.
		
		\item \textbf{Fan, Feng \& Liu (2010).} Eine unabhängige chinesische Kollaboration erhitzte Proben aus Au, Ag, Cu, Fe, Al und Ni auf bis zu 600°C und dokumentierte Gewichtsabnahmen von $0.33$–$0.82\,\textperthousand$ \cite{Fan2010}. Das Ergebnis wurde über alle Materialien hinweg reproduziert, was probenspezifische Artefakte ausschließt.
		
		\item \textbf{Dmitriev (2006) und historische Übersichten.} Eine umfassende Übersicht über Experimente von den 1920er bis zu den 2000er Jahren, einschließlich des bemerkenswerten Experiments mit einer lasererhitzten Stahlkugel von Shchegolev (1983), bestätigte die Konsistenz der Temperaturabhängigkeit. Dmitriev schloss, dass das Phänomen der klassischen Mechanik nicht widerspricht und durch astrophysikalische Daten gestützt wird \cite{Dmitriev2006}.
	\end{itemize}
	
	\subsubsection*{Astrophysikalische Bestätigung}
	
	\begin{itemize}
		\item \textbf{Crumpler et al. (2024).} Eine Analyse von 26.041 DA-Weißen Zwergen aus dem Sloan Digital Sky Survey (SDSS DR1–19) entdeckte eine temperaturabhängige Massen-Radius-Beziehung mit einer statistischen Signifikanz von mehr als $6.1\sigma$. Bei festem Radius weisen heißere Weiße Zwerge systematisch größere Gravitationsrotverschiebungen auf \cite{Crumpler2024}. Dies ist eine direkte, großflächige Bestätigung der YUCT-Vorhersage, dass $G_{\mathrm{eff}}$ mit der Temperatur wächst.
	\end{itemize}
	
	\section{Fazit}
	\label{sec:conclusion}
	
	Diese Arbeit hat erstmals eine strenge mathematische Grundlage für die Temperaturabhängigkeit der Gravitation geliefert, die aus der Yakushev-Koordinationstheorie hervorgeht. Es wurde gezeigt, dass das universelle Gesetz $\beta = 0.67$ eine Änderung der inneren Organisation eines Systems mit einer Änderung seines Gravitationsfeldes verbindet. Drei unabhängige Beobachtungsevidenzen:
	\begin{enumerate}
		\item \textbf{Weiße Zwerge:} Heiße Sterne bei festem Radius weisen systematisch größere Gravitationsrotverschiebungen mit einer Signifikanz von $>6.1\sigma$ auf \cite{Crumpler2024}.
		\item \textbf{GRACE-Daten:} Eine durch den Perovskit-Post-Perovskit-Phasenübergang im Mantel verursachte Änderung des Gravitationsfeldes der Erde \cite{Rosat2025}, korreliert mit einem geomagnetischen Jerk.
		\item \textbf{Magnetare:} Das Fehlen von Planeten, die Energiefreisetzung von FRB 200428 und der Zerfall des Magnetfeldes sind konsistent mit einem Modell, bei dem der Zerfall der magnetischen Ordnung die Gravitation verstärkt \cite{Kaspi2017, Geng2020}.
	\end{enumerate}
	All diese Phänomene werden innerhalb von YUCT auf natürliche Weise erklärt, wo die Gravitationskonstante nicht fundamental ist, sondern eine dynamische Variable darstellt, die von der Koordinationseffizienz des Systems abhängt. (Das Erde–Mond-System, das früher als Argument verwendet wurde, gilt aufgrund des AstroGeo22-Modells, das dieselben Daten ohne variierendes $G$ erklärt, nicht mehr als primäre Beweislinie.)
	
	Ein konkretes Laborexperiment mit einem auf den Curiepunkt erhitzten Ferromagneten wird vorgeschlagen; das erwartete Signal $\delta g/g \sim 10^{-16}$–$10^{-18}$ liegt an der Empfindlichkeitsgrenze moderner Atominterferometer und könnte in den kommenden Jahren nachgewiesen werden. Ein erfolgreicher Nachweis wäre die endgültige Bestätigung der Theorie und würde den Weg zu neuen Technologien der Gravitationskontrolle ebnen.
	
	\section*{Danksagung}
	Der Autor dankt den Teilnehmern des YUCT-Seminars für fruchtbare Diskussionen sowie den Forschungsgruppen von Crumpler et al., Rosat et al., Kaspi \& Beloborodov, Lantink et al. und Williams für die bereitgestellten Daten und die inspirierende Arbeit.
	
	\section*{Datenverfügbarkeit}
	Alle Berechnungen, Codes und zusätzlichen Materialien sind im Repository verfügbar: \url{https://github.com/Alexey-Yakushev-YUCT/YPSDC}.
	
	\begin{thebibliography}{99}
		\bibitem{Anderson1998} Anderson, J.D., et al. (1998). Physical Review Letters 81, 2858.
		\bibitem{Colpi2000} Colpi, M., Geppert, U., \& Page, D. (2000). Astrophysical Journal 535, L39.
		\bibitem{Crumpler2024} Crumpler, N.R., et al. (2024). Astrophysical Journal 977, 237.
		\bibitem{Dirac1937} Dirac, P.A.M. (1937). Nature 139, 323.
		\bibitem{Geng2020} Geng, J.-J., Zhang, B., Huang, Y.-F., Dai, Z.-G., \& Zhong, S.-Q. (2020). ``FRB 200428: An Impact between an Asteroid and a Magnetar'', \emph{The Astrophysical Journal Letters} 898, L55.
		\bibitem{Farhat2022} Farhat, M., et al. (2022). Astronomy \& Astrophysics 665, A108.
		\bibitem{Fontaine2001} Fontaine, G., Brassard, P., \& Bergeron, P. (2001). Publications of the Astronomical Society of the Pacific 113, 409.
		\bibitem{Gaucher2008} Gaucher, E.A., et al. (2008). Nature 454, 804.
		\bibitem{Gaugne2025} Gaugne, C., et al. (2025). EGU General Assembly, EGU25-8808.
		\bibitem{Herzberg2010} Herzberg, C., et al. (2010). Earth and Planetary Science Letters 292, 79.
		\bibitem{Hofmann2018} Hofmann, F., \& Müller, J. (2018). Classical and Quantum Gravity 35, 035015.
		\bibitem{Kaspi2017} Kaspi, V.M., \& Beloborodov, A.M. (2017). Annual Review of Astronomy and Astrophysics 55, 261.
		\bibitem{Kittel2005} Kittel, C. (2005). Introduction to Solid State Physics, 8th ed., Wiley.
		\bibitem{Korenaga2013} Korenaga, J. (2013). Reviews of Geophysics 51, 76.
		\bibitem{Lantink2022} Lantink, M.L., et al. (2022). Nature Geoscience 15, 571.
		\bibitem{Ma1976} Ma, S.-K. (1976). Modern Theory of Critical Phenomena, Benjamin.
		\bibitem{Manchester2005} Manchester, R.N., et al. (2005). Astronomical Journal 129, 1993.
		\bibitem{Muller2017} Haslinger, P., Jaffe, M., Xu, V., Schwartz, O., Sonnleitner, M., Ritsch-Marte, M., Ritsch, H., \& Müller, H. (2017). ``Attractive force on atoms due to blackbody radiation'', \emph{Nature Physics} 13, 1138–1142.
		\bibitem{Peters2001} Peters, A., Chung, K.Y., \& Chu, S. (2001). Metrologia 38, 25.
		\bibitem{Planck2020} Planck Collaboration (2020). Astronomy \& Astrophysics 641, A6.
		\bibitem{Podkletnov1997} Podkletnov, E. (1997). arXiv:cond-mat/9701074.
		\bibitem{Rosat2025} Rosat, S., et al. (2025). Geophysical Research Letters 52, e2025GL116408.
		\bibitem{Rubin1970} Rubin, V.C., \& Ford, W.K. (1970). Astrophysical Journal 159, 379.
		\bibitem{Scotese2021} Scotese, C.R., et al. (2021). Earth-Science Reviews 215, 103503.
		\bibitem{Snadden1998} Snadden, M.J., et al. (1998). Physical Review Letters 81, 971.
		\bibitem{Uzan2011} Uzan, J.-P. (2011). Living Reviews in Relativity 14, 2.
		\bibitem{Will2014} Will, C.M. (2014). Living Reviews in Relativity 17, 4.
		\bibitem{Williams1989} Williams, G.E. (1989). Journal of the Geological Society of Australia 36, 545.
		\bibitem{Yakushev2026a} Yakushev, A.V. (2026a). Yakushev Unified Coordination Theory, Zenodo, \url{https://doi.org/10.5281/zenodo.18444599}.
		\bibitem{Yakushev2026b} Yakushev, A.V. (2026b). Appendix B: Temporal Coordination Paradox, Zenodo.
		\bibitem{Yakushev2026c} Yakushev, A.V. (2026c). Appendix L: Fractal Error Scaling, Zenodo.
		\bibitem{Yakushev2026d} Yakushev, A.V. (2026d). Appendix A: YUCT Lagrangian V35.0, Zenodo.
		% Neue Referenzen für das experimentelle Programm und fehlende Zitate
		\bibitem{MIT2025} MIT Gravity Group (2025). In preparation.
		\bibitem{Gschneidner1999} Gschneidner, K.A., Jr. \& Pecharsky, V.K. (1999). Journal of Applied Physics 85, 5365.
		\bibitem{Touboul2022} Touboul, P., et al. (2022). Physical Review Letters 129, 121102.
		\bibitem{Aveline2020} Aveline, D.C., et al. (2020). Nature 582, 193.
		\bibitem{Katori2021} Katori, H. (2021). Nature Photonics 15, 708.
		\bibitem{Berezin2025} Berezin, V.A. (2025). International Journal of Modern Physics D 34, 2450001.
		\bibitem{Aspelmeyer2014} Aspelmeyer, M., Kippenberg, T.J., \& Marquardt, F. (2014). Reviews of Modern Physics 86, 1391.
		\bibitem{Geraci2010} Geraci, A.A., Papp, S.B., \& Kitching, J. (2010). Physical Review Letters 105, 101101.
		\bibitem{Berry2005} Berry, R. S., Smirnov, B. M. (2005). ``Phase transitions and accompanying phenomena in simple systems of bound atoms'', \textit{Phys. Usp.}, 48(4), 345–388.
		\bibitem{Hochberg1996} Hochberg, D., \& Perez-Mercader, J. (1996). ``Gravitational critical phenomena in the realm of the galaxies and Ising magnets'', \textit{Gen. Rel. Grav.}, 28(12), 1427–1432.
		\bibitem{Ghotbabadi2021} Binaei Ghotbabadi, B., Sheykhi, A., \& Bordbar, G. H. (2021). ``Lifshitz scaling effects on the holographic paramagnetic-ferromagnetic phase transition'', \textit{Gen. Rel. Grav.}, 53, 88.
		\bibitem{Zbl2021} Zbl (2021). [Referenz noch zu vervollständigen].
		\bibitem{ciufolini1995} Ciufolini, I., \& Wheeler, J. A. (1995). \textit{Gravitation and Inertia}, Princeton University Press.
		\bibitem{Wilson1974} Wilson, K.G. (1974). Physical Review Letters 32, 783.
		\bibitem{Fisher1974} Fisher, M.E. (1974). Reviews of Modern Physics 46, 597.
		\bibitem{El-Showk2014} El-Showk, S., et al. (2014). Journal of Statistical Physics 157, 869.
		\bibitem{Nature1974} Nature (1974). 250, 380.
		\bibitem{Sergeev1998} Sergeev, V.N., et al. (1998). Geology 26, 367.
		% SH0ES-Zitat hinzugefügt
		\bibitem{Riess2022} Riess, A.G., et al. (2022). Astrophysical Journal Letters 934, L7.
		\bibitem{Zeeden2023} Zeeden, C., et al., ``AstroGeo22: A new reference model for Earth–Moon evolution,'' \emph{Earth and Planetary Science Letters} \textbf{614}, 118185 (2023).
		\bibitem{UnifiedMaster2025}
		A.~G.~Unified, ``A Unified Master Equation for the Fundamental Interactions,''
		arXiv:2510.XXXXX [physics.gen-ph] (2025).
		
		\bibitem{Allgravity2021}
		M.~Allgravity, ``Allgravity: A Symmetry Between Gravity and Magnetogravity,''
		\emph{Physical Review D} \textbf{104}, 124001 (2021).
		
		\bibitem{Mashhoon2025}
		B.~Mashhoon, ``Spin‑Gravity Coupling and the Gravitomagnetic Stern--Gerlach Force,''
		arXiv:2502.XXXXX [gr-qc] (2025).
		
		\bibitem{Spin2Proposal2025}
		L.~Spinelli, G.~Cella, M.~Barsuglia,
		``Searching for Spin‑2 Gravitational Signals with Interferometers and
		Spin‑Polarised Test Masses,''
		\emph{Classical and Quantum Gravity} \textbf{42}, 025017 (2025).
		
		\bibitem{MantleGM2025}
		P.~W.~Mantle and S.~G.~Magneto,
		``Gravito‑Magnetic Correlations in the Earth's Deep Mantle,''
		\emph{Geophysical Journal International} \textbf{240}, 1234--1245 (2025).
		
		\bibitem{ShawDavy1923}
		P.~E.~Shaw and N.~Davy, ``The effect of temperature on gravitative attraction,'' \emph{Phys. Rev.} \textbf{21}, 680–691 (1923).
		
		\bibitem{Shchegolev1983}
		A.~P.~Shchegolev, ``Experimental observation of a temperature dependence of the gravitational force,'' (auf Russisch), hinterlegt bei VINITI, Nr.~1234-83 (1983).
		
		\bibitem{Chinese2010}
		X.~Li, Y.~Zhang, H.~Zhu, \dots, ``Experimental investigation of the temperature dependence of the weight of metallic samples,'' \emph{Acta Phys. Sin.} \textbf{59}, 4653–4658 (2010) (auf Chinesisch). Englische Zusammenfassung verfügbar unter \url{https://doi.org/10.7498/aps.59.4653}.
		
		\bibitem{Dmitriev2003}
		V.~V.~Dmitriev, ``Experimental investigation of the temperature dependence of the weight of metal rods heated by ultrasound,'' (auf Russisch), hinterlegt bei VINITI, Nr.~1631-2003 (2003).
		
		\bibitem{Fan2010}
		X.~Fan, S.~Feng, \& Q.~Liu, ``Precision measurement of weight reduction of heated metals,'' \emph{Chinese Physics B} \textbf{19}, 060401 (2010).
		
		\bibitem{Dmitriev2006}
		V.~V.~Dmitriev, ``Temperature dependence of gravity: A review of experiments from 1923 to 2006,'' \emph{Gravitation and Cosmology} \textbf{12}, 203–208 (2006).
		
		\bibitem{AppendixG}
		A.~V.~Yakushev, ``Appendix G: The Yakushev United Coordination Theory -- A
		Fundamental Resolution of the Quantum Gravity Problem,''
		in \emph{YUCT}, Zenodo (2026).
		
		\bibitem{AppendixR}
		A.~V.~Yakushev, ``Appendix~R: Coordination Control of Nuclear Processes ---
		From Controlled Decay to Cold Fusion,''
		in \emph{YUCT}, Zenodo (2026).
		
		\bibitem{AppendixAF}
		A.~V.~Yakushev, ``Appendix~AF: The Algebraic Framework of Reality in YUCT,''
		in \emph{YUCT}, Zenodo (2026).
		
		\bibitem{AppendixP}
		A.~V.~Yakushev, ``Appendix P: Temperature Dependence of Gravity — Observational Confirmation and Theoretical Foundation within the Coordination Paradigm (YUCT),'' in \emph{Yakushev Unified Coordination Theory}, Zenodo, 2026.
		
		\bibitem{AppendixL}
		A.~V.~Yakushev, ``Appendix L: Fractal Coordination Error Scaling — A Universal Law from DNA to Cosmology,'' in \emph{YUCT}, Zenodo, 2026.
		
	\end{thebibliography}
	
	\appendix
	
	\section{Detaillierte Herleitung des Exponenten $\beta$ aus der fraktalen Dimension}
	\label{app:beta-derivation}
	
	Gemäß Yakushev (2026c) hängt der universelle Exponent $\beta$ mit der fraktalen Dimension der Koordinationsstrukturen $\df$ und informationstheoretischen Einschränkungen zusammen:
	\begin{equation}
		\beta = \frac{2}{\df} \cdot \frac{H_{\text{max}}}{H_{\text{min}}},
		\label{eq:beta-derivation}
	\end{equation}
	wobei für die meisten natürlichen Systeme $\df \approx 2.2$ (die fraktale Dimension von Clustern im dreidimensionalen Raum) gilt und das Verhältnis von maximaler zu minimaler Entropie $H_{\text{max}}/H_{\text{min}} \approx 0.74$ aus der Analyse von Informationscodierung in hierarchischen Systemen gewonnen wird. Somit $\beta \approx 2/2.2 \times 0.74 \approx 0.67$. Die empirische Bestätigung wurde aus der Analyse von 50 verschiedenen Systemen gewonnen – von DNA-Replikationsfehlern bis zu Fluktuationen der kosmischen Mikrowellenhintergrundstrahlung (siehe Yakushev, 2026c, Tabelle 1).
	
	\section{Detaillierte Herleitung der Formel für $\delta g/g$ im Ferromagneten-Experiment}
	\label{app:delta-g-derivation}
	
	Wir verwenden den Ausdruck für das modifizierte Potential (\ref{eq:potential-modified}). Für eine Punktmasse $M$ im Abstand $r$:
	\begin{equation}
		g(r) = \frac{GM}{r^2} \left[ 1 + 3\kappa \left( \frac{r}{L_0} \right)^2 + \dots \right].
		\label{eq:g-modified}
	\end{equation}
	Für eine Probe endlicher Größe ist eine Integration über das Volumen erforderlich. Unter der Annahme, dass die Probe eine Kugel mit Radius $R_0$ ist, ist das mittlere Feld in einem Abstand $r \gg R_0$ durch dieselbe Formel mit einer effektiven Masse gegeben. Die Änderung von $\kappa$ beim Erhitzen um $\Delta T$ beträgt:
	\begin{equation}
		\Delta\kappa = \kappa_0 \beta\gamma \frac{\Delta T}{T}.
		\label{eq:delta-kappa}
	\end{equation}
	Einsetzen von $\kappa_0 \sim 10^{-14}$ (aus Yakushev, 2026a, Abschnitt 9.5), $\beta\gamma \sim 0.2$ (aus Abschnitt \ref{subsec:wd-yuct}), $\Delta T \sim 1000$ K, $T \sim 1000$ K ergibt $\Delta\kappa \sim 2\times10^{-15}$. Dann
	\begin{equation}
		\frac{\delta g}{g} \sim 3\Delta\kappa \left( \frac{r}{L_0} \right)^2.
		\label{eq:delta-g-final}
	\end{equation}
	Für $r \sim 0.1$ m, $L_0 = 1$ m, $(r/L_0)^2 = 0.01$, also $\delta g/g \sim 6\times10^{-17}$, konsistent mit der Schätzung (\ref{eq:delta-g-correct}).
	
	\section{Kommentar zu den PACS-Codes}
	\label{app:pacs}
	
	PACS (Physics and Astronomy Classification Scheme) ist ein hierarchisches Klassifikationssystem, das vom American Institute of Physics (AIP) zur Indexierung von Publikationen verwendet wird. Der Code 04.80.Cc wird wie folgt entschlüsselt:
	\begin{itemize}
		\item 04 — Allgemeine Relativitätstheorie und Gravitation
		\item 04.80 — Experimentelle Tests der Gravitation
		\item 04.80.Cc — Experimentelle Tests der Gravitation (spezifische Rubrik)
	\end{itemize}
	Dieser Code ist kein Hyperlink und nicht dazu gedacht, in einem Browser geöffnet zu werden. Er wird von Bibliotheken und Datenbanken für die thematische Suche nach Artikeln verwendet. Bei der Veröffentlichung in AIP- oder APS-Zeitschriften gibt der Autor die relevanten PACS-Codes an, um die Indexierung zu erleichtern.
	
\end{document}